% !TEX root = AImath.v4.tex
\chapter*{ 前言}
%\markboth{前言}{前言}

\addcontentsline{toc}{chapter}{前言}

 \begin{introduction}
  \item Definition of Theorem
  \item Ask for help
  \item Optimization Problem
  \item Property of Cauchy Series
  \item Angle of Corner
\end{introduction}

现在看起来的“显而易见”，在当年却是“困难重重”。


恰如其分地拿捏尺度，显然是智能的外在表现之一。“过犹不及”说得就是这个理。



“语言是思想的边界，技术是思想的实现。”


人工智能的历史源远流长。在古代的神话传说中，技艺高超的工匠可以制作人造“人（robot）”，并为其赋予智能或意识。



本质上，人类发展史就是一部智能简史，人类从计数开始，到数学理论的发展，都在不断追求计算越来越高效、越来越自动化，越来越智能...。从人类文明最早出现数，数的标记和运算，从算术到代数、几何、抽象代数、拓扑、群论等越来越复杂和专门的数学分支，数学史的发展主线背后，是对智能的极致追求。



从人脑和手算发展到一定程度，对更高效计算效率的追求必然会导致计算机的发明，导致人工智能



这个思路非常有启发性，您将人类发展史看作是一部智能简史，从最早的数和运算到现代计算机和人工智能的发展，强调了数学和智能追求之间的紧密联系。我们可以将这个思路融入到书的开篇和部分章节中，甚至作为书的核心主题来展开讨论。以下是对目录的调整建议，结合您的思路：



这个思路非常有启发性，您将人类发展史看作是一部智能简史，从最早的数和运算到现代计算机和人工智能的发展，强调了数学和智能追求之间的紧密联系。我们可以将这个思路融入到书的开篇和部分章节中，甚至作为书的核心主题来展开讨论。以下是对目录的调整建议，结合您的思路：




-它们在数百年的黑暗中闪耀，用数学照亮了人类的历程。

在伏尔泰的小说《老实人》中，坚持莱布尼兹式乐观主义的邦葛罗斯安慰幸存者时，说“如果没有这场灾难，事物就无法继续进行了。因为一切都是最好的安排。”

歌德把历史称作“上帝的神秘作坊”。在这个“作坊”里，堆积这对人类来说无关紧要的小事。只在极少时候，它们才会被某些时刻着凉，史蒂芬·茨威格称之为“高光时刻”。

每一个“宿命时刻”，都是一个个事件串联的了解过。一连串的事件被紧密地编织在一起，经过漫长的发展，事件一个接着一个，然后突然成熟。历史这个“神秘作坊”不只生成武器和战争。其中一些“宿命时刻”如繁星办“熠熠生辉、永不熄灭”。

这本书代表了我们的尝试——让深度学习可平易近人，教会人们概念、背景和代码。

一种结合了代码、数学和HTML的媒介

任何一种计算技术要想发挥其全部影响力，都必须得到充分的理解、充分的文档记录，并得到成熟的、维护良好的工具的支持。关键思想应该被清楚地提炼出来，尽可能减少需要让新的从业者跟上时代的入门时间。成熟的库应该自动化常见的任务，示例代码应该使从业者可以轻松地修改、应用和扩展常见的应用程序，以满足他们的需求。以动态网页应用为例。尽管许多公司，如亚马逊，在20世纪90年代开发了成功的数据库驱动网页应用程序。但在过去的10年里，这项技术在帮助创造性企业家方面的潜力已经得到了更大程度的发挥，部分原因是开发了功能强大、文档完整的框架。

测试深度学习的潜力带来了独特的挑战，因为任何一个应用都会将不同的学科结合在一起。应用深度学习需要同时了解（1）以特定方式提出问题的动机；（2）给定建模方法的数学；（3）将模型拟合数据的优化算法；（4）能够有效训练模型、克服数值计算缺陷并最大限度地利用现有硬件的工程方法。同时教授表述问题所需的批判性思维技能、解决问题所需的数学知识，以及实现这些解决方案所需的软件工具，这是一个巨大的挑战。


-数学思维与人工智能的关系

-从数学到智能的发展历程

\section{内容和结构}

全书大致可分为三个部分，在[图1](https：//zh-v2.d2l.ai/chapter\_preface/index.html fig-book-org)中用不同的颜色呈现：

%![../_images/book-org.svg](https：//zh-v2.d2l.ai/_images/book-org.svg)

图1全书结构

-第一部分包括基础知识和预备知识。[1节]提供深度学习的入门课程。然后在[2节]中，我们将快速介绍实践深度学习所需的前提条件，例如如何存储和处理数据，以及如何应用基于线性代数、微积分和概率基本概念的各种数值运算。[3节] 和[4节]涵盖了深度学习的最基本概念和技术，例如线性回归、多层感知机和正则化。

-接下来的五章集中讨论现代深度学习技术。[5节]描述了深度学习计算的各种关键组件，并为我们随后实现更复杂的模型奠定了基础。接下来，在[6节]和[7节]中，我们介绍了卷积神经网络（convolutionalneuralnetwork，CNN），这是构成大多数现代计算机视觉系统骨干的强大工具。随后，在[8节]和[9节]中，我们引入了循环神经网络(recurrentneuralnetwork，RNN)，这是一种利用数据中的时间或序列结构的模型，通常用于自然语言处理和时间序列预测。在[10节]中，我们介绍了一类新的模型，它采用了一种称为注意力机制的技术，最近它们已经开始在自然语言处理中取代循环神经网络。这一部分将帮助读者快速了解大多数现代深度学习应用背后的基本工具。

-第三部分讨论可伸缩性、效率和应用程序。首先，在[11节]中，我们讨论了用于训练深度学习模型的几种常用优化算法。下一章[12节]将探讨影响深度学习代码计算性能的几个关键因素。在[13节]中，我们展示了深度学习在计算机视觉中的主要应用。在[14节]和[15节]中，我们展示了如何预训练语言表示模型并将其应用于自然语言处理任务。

     小结
-深度学习已经彻底改变了模式识别，引入了一系列技术，包括计算机视觉、自然语言处理、自动语音识别。

-要成功地应用深度学习，必须知道如何抛出一个问题、建模的数学方法、将模型与数据拟合的算法，以及实现所有这些的工程技术。

-这本书提供了一个全面的资源，包括文本、图表、数学和代码，都集中在一个地方。

-要回答与本书相关的问题，请访问我们的论坛[discuss.d2l.ai](https：//discuss.d2l.ai/).

-所有Jupyter记事本都可以在GitHub上下载。



-本书目的与目标读者

-数学思维与人工智能的关系

-从数学到智能的发展历程

    第一部分：人类智能简史

2.第1章：智能简史

-1.2数学的发展与智能的进化
-表示虚无的数字0

“印度数学家的成果跟随阿拉伯人在地中海沿岸传播了好久，才最终来到中世纪欧洲翻译家的手中。穆罕默德·花拉子米是一位数学家，出生于乌兹别克斯坦，生活于780年至850年间，著有《代数学》（Al-kitābal-mukhtaṣarfīḥisābal-ğabrwa’l-muqābala）。代数（algebra）一词出自书名中的阿拉伯文“al-jabr”的拉丁转写，演算法（Algorism）则出自作者花拉子米（al-Khwārizmī）的拉丁文译名。在《代数学》这部论著中，没有公式，没有零，甚至连九个印度数字也没有出现。而在公元8世纪以前，阿拉伯人使用的是从希腊人那里学来的记数法，传统的希腊数字是用字母表示的。花拉子米在另一部作品——《印度的计算术》中提到了零，“圆圈即无”。我们如今只能看到这部作品的拉丁文版本，与遗失的阿拉伯文真本相比，它在内容上做了相当多的改动。”

“人们何时发现正方形或长方形的对角线与边长之比无法用整数表示，“无限”随之闯入数学的世界？
这个比例为何又出现在了由一个兔子问题而衍生出的一串数字中，经斐波那契之手成了黄金比例？”

-发现无理数

-一桩逻辑上的丑闻

-虚数

-一个关于秘密和背叛的故事

-把方程式变成诗

-计算虚数

-非欧几何的世界

-不止于三维

-“真正”的空间几何

-“想象“几何

-弯曲空间

-抽象数学与智能的拓展

-1.3从手算到计算机

-计算工具的进化

-对高效计算的追求如何促成计算机的发明

-1.4人工智能的诞生

-从自动化到智能化的跨越

-人工智能的早期发展与数学的支撑

    第二部分：人工智能中的核心数学思维

1.第2章：抽象与建模：从现实到数学

-2.1建模思维：如何从现实问题抽象出数学模型

-2.2常见的AI模型及其数学思维：如线性模型、非线性模型等

-2.3模型的简化与优化：模型复杂度与可解释性的平衡

2.第3章：数据的表示与处理：几何与代数的视角

-3.1数据表示与向量空间：如何通过几何思维理解数据

-3.2线性代数在AI中的角色：从矩阵运算到特征提取

-3.3特征值、特征向量与数据降维：PCA、SVD等技术背后的数学思维

3.第4章：不确定性与决策：概率与统计的框架

-4.1概率思维：处理不确定性的数学工具

-4.2统计推断：如何从数据中学习

4.第5章：优化与决策思维：在复杂环境下做出最佳决策

-最优化问题的数学框架

-凸优化与非凸优化的区别

-启发式算法与实用优化方法

5.第6章：学习与推理：从算法到智能

-5.1机器学习中的数学思维：监督学习与无监督学习的数学基础

-5.2推理与知识表示：逻辑、概率与其他形式的推理方法

-5.3从理论到实践：算法的优化与应用

6.第7章：算法设计中的计算思维

-经典AI算法的数学分析

-计算复杂度：NP问题与算法效率

-计算思维如何影响算法设计

7.第8章：深度学习中的数学直觉

-神经网络的数学基础：从线性代数到非线性激活

-深度学习的几何与拓扑视角：卷积、池化与拓扑思维

-深度学习的优化挑战：反向传播与梯度下降的数学理解

    第三部分：应用与与案例分析

1.第9章：AI中的复杂性与可解释性

-可解释性问题：从数学角度审视AI的可解释性

-实例分析：如何在实际应用中平衡复杂性与可解释性

2.第10章：数据处理与统计推断

-数据分析中的数学思维

-统计推断与机器学习模型

-从数据到决策的思维路径

    第四部分：未来展望与思考
    
10.第11章：智能的未来与数学思维

-数学思维如何推动智能的发展

-智能的前沿挑战：AI发展中的数学难题

-从数学思维到未来创新：如何利用数学思维推动AI的下一步发展

11.第12章：结语

-从数学思维到人工智能的未来

-如何在实践中培养数学思维




在这个框架中，第一部分特别突出了人类智能发展的历史背景，并将其与数学的发展相结合，形成一个贯穿全书的主线。这不仅帮助读者理解人工智能的本质，还能引导他们思考数学在智能追求中的关键作用。您觉得这个调整是否符合您的思路呢？
===

中道：

势道术 or 道法术器势：道以明向，法以立本，术以立策，器以成事，势以立人


应用深度学习需要同时了解

（1）以特定方式提出问题的动机；

（2）给定建模方法的数学；

（3）将模型拟合数据的优化算法；

（4）能够有效训练模型、克服数值计算缺陷并最大限度地利用现有硬件的工程方法。

要成功地应用深度学习，必须知道如何抛出一个问题、建模的数学方法、将模型与数据拟合的算法，以及实现所有这些的工程技术。

这个过程包含了数学思维、算法思维、工程思维。



或者只有教科书，或者只有代码没有数学原理

我们发现了大量关于如何使用给定的深度学习框架（例如，如何对TensorFlow中的矩阵进行基本的数值计算)或实现特定技术的代码示例（例如，LeNet、AlexNet、ResNet的代码片段），这些代码示例分散在各种博客帖子和GitHub库中。

但是，这些示例通常关注如何实现给定的方法，但忽略了为什么做出某些算法决策的讨论。

===

就在几年前，不管在大公司还是创业公司，都鲜有工程师和科学家来将深度学习应用到智能产品与服务中。作为深度学习前身的神经网络，才刚刚摆脱被机器学习学术界认为是过时工具的印象。那个时候，即使是机器学习也非新闻头条的常客。它仅仅被看作是一门具有前瞻性，并拥有一系列小范围实际应用的学科。在包含计算机视觉和自然语言处理在内的实际应用通常需要大量的相关领域知识：这些实际应用被视为相互独立的领域，而机器学习只占其中一小部分。

然而仅仅在这几年之内，深度学习便令全世界大吃一惊。它非常有力地推动了计算机视觉、自然语言处理、自动语音识别、强化学习和统计建模等多个领域的快速发展。随着这些领域的不断进步，我们现在可以制造自动驾驶的汽车，基于短信、邮件甚至电话的自动回复系统，以及在围棋中击败最优秀人类选手的软件。这些由深度学习带来的新工具也正产生着广泛的影响：它们改变了电影制作和疾病诊断的方式，并在从天体物理学到生物学等各个基础科学中扮演越来越重要的角色。

与此同时，深度学习也给它的使用者们带来了独一无二的挑战：任何单一的应用都汇集了各学科的知识。具体来说，应用深度学习需要同时理解：

1. 问题的动机和特点；
2. 将大量不同类型神经网络层通过特定方式组合在一起的模型背后的数学原理；
3. 在原始数据上拟合极复杂的深层模型的优化算法；
4. 有效训练模型、避免数值计算陷阱以及充分利用硬件性能所需的工程技能；
5. 为解决方案挑选合适的变量（超参数）组合的经验。

同样，我们几位作者也面临前所未有的挑战：我们需要在有限的篇幅里糅合深度学习的多方面知识，从而使读者能够较快理解并应用深度学习技术。本书代表了我们的一种尝试：我们将教给读者概念、背景知识和代码；我们将在同一个地方阐述剖析问题所需的批判性思维、解决问题所需的数学知识，以及实现解决方案所需的工程技能。

 包含代码、数学、网页、讨论的统一资源

我们在2017年7月启动了写作这本书的项目。当时我们需要向用户解释ApacheMXNet在那时的新接口Gluon。遗憾的是，我们并没有找到任何一个资源可以同时满足以下几点需求：

1. 包含较新的方法和应用，并不断更新；
2. 广泛覆盖现代深度学习技术并具有一定的技术深度；
3. 既是严谨的教科书，又是包含可运行代码的生动的教程。

那时，我们在博客和GitHub上找到了大量的演示特定深度学习框架（例如用TensorFlow进行数值计算）或实现特定模型（例如AlexNet、ResNet等）的示例代码。这些示例代码的一大价值在于提供了教科书或论文往往省略的实现细节，比如数据的处理和运算的高效率实现。如果不了解这些，即使能将算法倒背如流，也难以将算法应用到自己的项目中去。此外，这些示例代码还使得用户能通过观察修改代码所导致的结果变化而快速验证想法、积累经验。因此，我们坚信动手实践对于学习深度学习的重要性。然而可惜的是，这些示例代码通常侧重于如何实现给定的方法，却忽略了有关算法设计的探究或者实现细节的解释。虽然在像Distill这样的网站和某些博客上出现了一些有关算法设计和实现细节的讨论，但它们常常缺少示例代码，并通常仅覆盖深度学习的一小部分。

另外，我们欣喜地看到了一些有关深度学习的教科书不断问世，其中最著名的要数Goodfellow、Bengio和Courville的《深度学习》。该书梳理了深度学习背后的众多概念与方法，是一本极为优秀的教材。然而，这类资源并没有将概念描述与实际代码相结合，以至于有时会令读者对如何实现它们感到毫无头绪。除了这些以外，商业课程提供者们虽然制作了众多的优质资源，但它们的付费门槛依然令不少用户望而生畏。

正因为这样，深度学习用户，尤其是初学者，往往不得不参考来源不同的多种资料。例如，通过教科书或者论文来掌握算法及其相关数学知识，阅读线上文档学习深度学习框架的使用方法，然后寻找感兴趣的算法在这个框架上的实现并摸索如何将它应用到自己的项目中去。如果你正亲身经历这一过程，你可能会感到痛苦：不同来源的资料有时难以相互一一对应，即便能够对应也可能需要花费大量的精力。例如，我们需要将某篇论文公式中的数学变量与某段网上实现中的程序变量一一对应，并在代码中找到论文可能没交代清楚的实现细节，甚至要为运行不同的代码安装不同的运行环境。

针对以上存在的痛点，我们正在着手创建一个为实现以下目标的统一资源：

1. 所有人均可在网上免费获取；
2. 提供足够的技术深度，从而帮助读者实际成为深度学习应用科学家：既理解数学原理，又能够实现并不断改进方法；
3. 包含可运行的代码，为读者展示如何在实际中解决问题。这样不仅直接将数学公式对应成实际代码，而且可以修改代码、观察结果并及时获取经验；
4. 允许我们和整个社区不断快速迭代内容，从而紧跟仍在高速发展的深度学习领域；
5. 由包含有关技术细节问答的论坛作为补充，使大家可以相互答疑并交换经验。

这些目标往往互有冲突：公式、定理和引用最容易通过LaTeX进行管理和展示，代码自然应该用简单易懂的Python描述，而网页本身应该是一堆HTML及配套的CSS和JavaScript。此外，我们希望这个资源可以作为可执行代码、实体书以及网站。然而，目前并没有任何工具可以完美地满足以上所有需求。

因此，我们不得不自己来集成这样的一个工作流。我们决定在GitHub上分享源代码并允许提交编辑，通过Jupyter记事本来整合代码、公式、文本、图片等，使用Sphinx作为渲染引擎来生成不同格式的输出，并使用Discourse作为论坛。虽然我们的系统尚未完善，但这些选择在互有冲突的目标之间取得了较好的折中。这很可能是使用这种集成工作流发布的第一本书。

 从在线课程到纸质书

本书的两位中国作者曾每周末在线免费讲授“动手学深度学习”系列课程。课程的讲义自然成为了本书内容的蓝本。这个课程持续了5个月，其间近3,000名同学参与了讨论，并贡献了近5,000多个有价值的讨论，特别是其中几个参加比赛的练习很受欢迎。这个课程的受欢迎程度出乎我们的意料。尽管我们将课件和课程视频都公开在了网上，但我们同时觉得出版成纸质书也许能让更多喜爱纸质阅读的读者受益。因此，我们委托人民邮电出版社来出版这本书。

从蓝本到成书花费了更多的时间。我们对所有涉及的所有技术点补充了背景介绍，并使用了更加严谨的写作风格，还对版式和示意图做了大量修改。书中所有的代码执行结果都是自动生成的，任何改动都会触发对书中每一段代码的测试，以保证读者在动手实践时能复现结果。

我们的初衷是让更多人更容易地使用深度学习。为了让大家能够便利地获取这些资源，我们保留了免费的网站内容，并且通过不收取出版稿费的方式来降低纸质书的价格，使更多人有能力购买。

 

 教学资源、计算资源和反馈

本书的英文版Dive into Deep Learning是加州大学伯克利分校2019年春学期“Introduction to Deep Learning”（深度学习导论）课程的教材。同时，截至2019年春学期，本书中英文版的内容已被全球15所知名大学用于教学。本书的学习社区、免费教学资源（课件、教学视频、更多习题等），以及用于本书学习或教学的免费计算资源（仅限学生和老师）的申请方法在本书网站 [https://zh.d2l.ai](https://zh.d2l.ai/) 上发布。诚然，将算法、公式、图片、代码和样例统一进一本适合阅读的书，并以具有交互式体验的Jupyter记事本文件的形式提供给读者，是对我们的极大挑战。书中难免有很多疏忽的地方，敬请原谅，并希望读者能通过每一节后面的二维码向我们反馈阅读本书过程中发现的问题。

结尾处，附上陆游的一句诗作为勉励：

“纸上得来终觉浅，绝知此事要躬行。”

===

 1. 引言

中道：



势道术 or 道法术器势：道以明向，法以立本，术以立策，器以成事，势以立人





第二版

 前言

几年前，在大公司和初创公司中，并没有大量的深度学习科学家开发智能产品和服务。我们中年轻人（作者）进入这个领域时，机器学习并没有在报纸上获得头条新闻。我们的父母根本不知道什么是机器学习，更不用说为什么我们可能更喜欢机器学习，而不是从事医学或法律职业。机器学习是一门具有前瞻性的学科，在现实世界的应用范围很窄。而那些应用，例如语音识别和计算机视觉，需要大量的领域知识，以至于它们通常被认为是完全独立的领域，而机器学习对这些领域来说只是一个小组件。因此，神经网络——我们在本书中关注的深度学习模型的前身，被认为是过时的工具。

就在过去的五年里，深度学习给世界带来了惊喜，推动了计算机视觉、自然语言处理、自动语音识别、强化学习和统计建模等领域的快速发展。有了这些进步，我们现在可以制造比以往任何时候都更自主的汽车（不过可能没有一些公司试图让大家相信的那么自主），可以自动起草普通邮件的智能回复系统，帮助人们从令人压抑的大收件箱中解放出来。在围棋等棋类游戏中，软件超越了世界上最优秀的人，这曾被认为是几十年后的事。这些工具已经对工业和社会产生了越来越广泛的影响，改变了电影的制作方式、疾病的诊断方式，并在基础科学中扮演着越来越重要的角色——从天体物理学到生物学。

 关于本书

这本书代表了我们的尝试——让深度学习可平易近人，教会人们概念、背景和代码。

 一种结合了代码、数学和HTML的媒介

任何一种计算技术要想发挥其全部影响力，都必须得到充分的理解、充分的文档记录，并得到成熟的、维护良好的工具的支持。关键思想应该被清楚地提炼出来，尽可能减少需要让新的从业者跟上时代的入门时间。成熟的库应该自动化常见的任务，示例代码应该使从业者可以轻松地修改、应用和扩展常见的应用程序，以满足他们的需求。以动态网页应用为例。尽管许多公司，如亚马逊，在20世纪90年代开发了成功的数据库驱动网页应用程序。但在过去的10年里，这项技术在帮助创造性企业家方面的潜力已经得到了更大程度的发挥，部分原因是开发了功能强大、文档完整的框架。

测试深度学习的潜力带来了独特的挑战，因为任何一个应用都会将不同的学科结合在一起。应用深度学习需要同时了解（1）以特定方式提出问题的动机；（2）给定建模方法的数学；（3）将模型拟合数据的优化算法；（4）能够有效训练模型、克服数值计算缺陷并最大限度地利用现有硬件的工程方法。同时教授表述问题所需的批判性思维技能、解决问题所需的数学知识，以及实现这些解决方案所需的软件工具，这是一个巨大的挑战。

在我们开始写这本书的时候，没有资源能够同时满足一些条件：（1）是最新的；（2）涵盖了现代机器学习的所有领域，技术深度丰富；（3）在一本引人入胜的教科书中，人们可以在实践教程中找到干净的可运行代码，并从中穿插高质量的阐述。我们发现了大量关于如何使用给定的深度学习框架（例如，如何对TensorFlow中的矩阵进行基本的数值计算)或实现特定技术的代码示例（例如，LeNet、AlexNet、ResNet的代码片段），这些代码示例分散在各种博客帖子和GitHub库中。但是，这些示例通常关注如何实现给定的方法，但忽略了为什么做出某些算法决策的讨论。虽然一些互动资源已经零星地出现以解决特定主题。例如，在网站[Distill](http://distill.pub/)上发布的引人入胜的博客帖子或个人博客，但它们仅覆盖深度学习中的选定主题，并且通常缺乏相关代码。另一方面，虽然已经出现了几本教科书，其中最著名的是 ([Goodfellow et al., 2016]
(https://zh-v2.d2l.ai/chapter\_references/zreferences.htmlid47))
（中文名《深度学习》），它对深度学习背后的概念进行了全面的调查，但这些资源并没有将这些概念的描述与这些概念的代码实现结合起来。有时会让读者对如何实现它们一无所知。此外，太多的资源隐藏在商业课程提供商的付费壁垒后面。

我们着手创建的资源可以：（1）每个人都可以免费获得；（2）提供足够的技术深度，为真正成为一名应用机器学习科学家提供起步；（3）包括可运行的代码，向读者展示如何解决实践中的问题；（4）允许我们和社区的快速更新;（5）由一个[论坛]
(http://discuss.d2l.ai/)
作为补充，用于技术细节的互动讨论和回答问题。

这些目标经常是相互冲突的。公式、定理和引用最好用LaTeX来管理和布局。代码最好用Python描述。网页原生是HTML和JavaScript的。此外，我们希望内容既可以作为可执行代码访问、作为纸质书访问，作为可下载的PDF访问，也可以作为网站在互联网上访问。目前还没有完全适合这些需求的工具和工作流程，所以我们不得不自行组装。我们在 [16.5节]中详细描述了我们的方法。我们选择GitHub来共享源代码并允许编辑，选择Jupyter记事本来混合代码、公式和文本，选择Sphinx作为渲染引擎来生成多个输出，并为论坛提供讨论。虽然我们的体系尚不完善，但这些选择在相互冲突的问题之间提供了一个很好的妥协。我们相信，这可能是第一本使用这种集成工作流程出版的书。

 在实践中学习

许多教科书教授一系列的主题，每一个都非常详细。例如，Chris Bishop的优秀教科书 ([Bishop, 2006]
(https://zh-v2.d2l.ai/chapter\_references/zreferences.htmlid10)) 
，对每个主题都教得很透彻，以至于要读到线性回归这一章需要大量的工作。虽然专家们喜欢这本书正是因为它的透彻性，但对初学者来说，这一特性限制了它作为介绍性文本的实用性。

在这本书中，我们将适时教授大部分概念。换句话说，你将在实现某些实际目的所需的非常时刻学习概念。虽然我们在开始时花了一些时间来教授基础的背景知识，如线性代数和概率，但我们希望你在思考更深奥的概率分布之前，先体会一下训练模型的满足感。

除了提供基本数学背景速成课程的几节初步课程外，后续的每一章都介绍了适量的新概念，并提供可独立工作的例子——使用真实的数据集。这带来了组织上的挑战。某些模型可能在逻辑上组合在单节中。而一些想法可能最好是通过连续允许几个模型来传授。另一方面，坚持“一个工作例子一节”的策略有一个很大的好处：这使你可以通过利用我们的代码尽可能轻松地启动你自己的研究项目。只需复制这一节的内容并开始修改即可。

我们将根据需要将可运行代码与背景材料交错。通常，在充分解释工具之前，我们常常会在提供工具这一方面犯错误（我们将在稍后解释背景）。例如，在充分解释随机梯度下降为什么有用或为什么有效之前，我们可以使用它。这有助于给从业者提供快速解决问题所需的弹药，同时需要读者相信我们的一些决定。

这本书将从头开始教授深度学习的概念。有时，我们想深入研究模型的细节，这些的细节通常会被深度学习框架的高级抽象隐藏起来。特别是在基础教程中，我们希望读者了解在给定层或优化器中发生的一切。在这些情况下，我们通常会提供两个版本的示例：一个是我们从零开始实现一切，仅依赖张量操作和自动微分；另一个是更实际的示例，我们使用深度学习框架的高级API编写简洁的代码。一旦我们教了您一些组件是如何工作的，我们就可以在随后的教程中使用高级API了。

 内容和结构

全书大致可分为三个部分，在 [图1]
(https://zh-v2.d2l.ai/chapter\_preface/index.htmlfig-book-org) 
中用不同的颜色呈现：
%![../\_images/book-org.svg](https://zh-v2.d2l.ai/\_images/book-org.svg)

图1 全书结构

- 第一部分包括基础知识和预备知识。 [1节]
%(https://zh-v2.d2l.ai/chapter\_introduction/index.htmlchap-introduction) 
提供深度学习的入门课程。然后在 [2节]中，我们将快速介绍实践深度学习所需的前提条件，例如如何存储和处理数据，以及如何应用基于线性代数、微积分和概率基本概念的各种数值运算。 [3节]和 [4节]
涵盖了深度学习的最基本概念和技术，例如线性回归、多层感知机和正则化。
- 接下来的五章集中讨论现代深度学习技术。 [5节]
描述了深度学习计算的各种关键组件，并为我们随后实现更复杂的模型奠定了基础。接下来，在 [6节]
和 [7节]
中，我们介绍了卷积神经网络（convolutional neural network，CNN），这是构成大多数现代计算机视觉系统骨干的强大工具。随后，在 [8节]和 [9节]
中，我们引入了循环神经网络(recurrent neural network，RNN)，这是一种利用数据中的时间或序列结构的模型，通常用于自然语言处理和时间序列预测。在 [10节]
中，我们介绍了一类新的模型，它采用了一种称为注意力机制的技术，最近它们已经开始在自然语言处理中取代循环神经网络。这一部分将帮助读者快速了解大多数现代深度学习应用背后的基本工具。
- 第三部分讨论可伸缩性、效率和应用程序。 首先，在 [11节]
中，我们讨论了用于训练深度学习模型的几种常用优化算法。下一章 [12节]
将探讨影响深度学习代码计算性能的几个关键因素。在 [13节]
中，我们展示了深度学习在计算机视觉中的主要应用。在 [14节]
和 [15节] 中，我们展示了如何预训练语言表示模型并将其应用于自然语言处理任务。



 代码

本书的大部分章节都以可执行代码为特色，因为我们相信交互式学习体验在深度学习中的重要性。
目前，某些直觉只能通过试错、小幅调整代码并观察结果来发展。
理想情况下，一个优雅的数学理论可能会精确地告诉我们如何调整代码以达到期望的结果。
不幸的是，这种优雅的理论目前还没有出现。
尽管我们尽了最大努力，但仍然缺乏对各种技术的正式解释，
这既是因为描述这些模型的数学可能非常困难，也是因为对这些主题的认真研究最近才进入高潮。
我们希望随着深度学习理论的发展，这本书的未来版本将能够在当前版本无法提供的地方提供见解。

有时，为了避免不必要的重复，我们将本书中经常导入和引用的函数、类等封装在`d2l`包中。
对于要保存到包中的任何代码块，比如一个函数、一个类或者多个导入，我们都会标记为`@save`。我们在 [16.6节]
(https://zh-v2.d2l.ai/chapter\_appendix-tools-for-deep-learning/d2l.htmlsec-d2l) 
中提供了这些函数和类的详细描述。`d2l`软件包是轻量级的，仅需要以下软件包和模块作为依赖项：

%[MXNET](https://zh-v2.d2l.ai/chapter\_preface/index.htmlmxnet-1-0)[PYTORCH](https://zh-v2.d2l.ai/chapter\_preface/index.htmlpytorch-1-1)[TENSORFLOW](https://zh-v2.d2l.ai/chapter\_preface/index.htmltensorflow-1-2)[PADDLE](https://zh-v2.d2l.ai/chapter\_preface/index.htmlpaddle-1-3)

```
@save
import collections
import hashlib
import math
import os
import random
import re
import shutil
import sys
import tarfile
import time
import zipfile
from collections import defaultdict
import pandas as pd
import requests
from IPython import display
%from matplotlib import pyplot as plt
%from matplotlib\_inline import backend\_inline

%d2l = sys.modules[\_\_name\_\_]
```

%![Copy to clipboard](https://raw.githubusercontent.com/choldgraf/sphinx-copybutton/master/sphinx\_copybutton/\_static/copy-button.svg)

 目标受众

本书面向学生（本科生或研究生）、工程师和研究人员，他们希望扎实掌握深度学习的实用技术。因为我们从头开始解释每个概念，所以不需要过往的深度学习或机器学习背景。全面解释深度学习的方法需要一些数学和编程，但我们只假设读者了解一些基础知识，包括线性代数、微积分、概率和非常基础的Python编程。此外，在附录中，我们提供了本书所涵盖的大多数数学知识的复习。大多数时候，我们会优先考虑直觉和想法，而不是数学的严谨性。有许多很棒的书可以引导感兴趣的读者走得更远。Bela Bollobas的《线性分析》 ([Bollobás, 1999]
(https://zh-v2.d2l.ai/chapter\_references/zreferences.htmlid13)) 
对线性代数和函数分析进行了深入的研究。 ([Wasserman, 2013]
(https://zh-v2.d2l.ai/chapter\_references/zreferences.htmlid178))
是一本很好的统计学指南。如果读者以前没有使用过Python语言，那么可以仔细阅读这个[Python教程]
%(http://learnpython.org/)。

 论坛

与本书相关，我们已经启动了一个论坛，在[discuss.d2l.ai](https://discuss.d2l.ai/)
。当对本书的任何一节有疑问时，请在每一节的末尾找到相关的讨论页链接。

 小结

- 深度学习已经彻底改变了模式识别，引入了一系列技术，包括计算机视觉、自然语言处理、自动语音识别。
- 要成功地应用深度学习，必须知道如何抛出一个问题、建模的数学方法、将模型与数据拟合的算法，以及实现所有这些的工程技术。
- 这本书提供了一个全面的资源，包括文本、图表、数学和代码，都集中在一个地方。
- 要回答与本书相关的问题，请访问我们的论坛。
- 所有Jupyter记事本都可以在GitHub上下载。


第一版

几年前，在大公司和初创公司中，并没有大量的深度学习科学家开发智能产品和服务。我们中年轻人（作者）进入这个领域时，机器学习并没有在报纸上获得头条新闻。我们的父母根本不知道什么是机器学习，更不用说为什么我们可能更喜欢机器学习，而不是从事医学或法律职业。机器学习是一门具有前瞻性的学科，在现实世界的应用范围很窄。而那些应用，例如语音识别和计算机视觉，需要大量的领域知识，以至于它们通常被认为是完全独立的领域，而机器学习对这些领域来说只是一个小组件。因此，神经网络——我们在本书中关注的深度学习模型的前身，被认为是过时的工具。

就在过去的五年里，深度学习给世界带来了惊喜，推动了计算机视觉、自然语言处理、自动语音识别、强化学习和统计建模等领域的快速发展。有了这些进步，我们现在可以制造比以往任何时候都更自主的汽车（不过可能没有一些公司试图让大家相信的那么自主），可以自动起草普通邮件的智能回复系统，帮助人们从令人压抑的大收件箱中解放出来。在围棋等棋类游戏中，软件超越了世界上最优秀的人，这曾被认为是几十年后的事。这些工具已经对工业和社会产生了越来越广泛的影响，改变了电影的制作方式、疾病的诊断方式，并在基础科学中扮演着越来越重要的角色——从天体物理学到生物学。

时至今日，人们常用的计算机程序几乎都是软件开发人员从零编写的。 比如，现在开发人员要编写一个程序来管理网上商城。 经过思考，开发人员可能提出如下一个解决方案： 首先，用户通过Web浏览器（或移动应用程序）与应用程序进行交互； 紧接着，应用程序与数据库引擎进行交互，以保存交易历史记录并跟踪每个用户的动态； 其中，这个应用程序的核心——“业务逻辑”，详细说明了应用程序在各种情况下进行的操作。

为了完善业务逻辑，开发人员必须细致地考虑应用程序所有可能遇到的边界情况，并为这些边界情况设计合适的规则。 当买家单击将商品添加到购物车时，应用程序会向购物车数据库表中添加一个条目，将该用户ID与商品ID关联起来。 虽然一次编写出完美应用程序的可能性微乎其微，但在大多数情况下，开发人员可以从上述的业务逻辑出发，编写出符合业务逻辑的应用程序，并不断测试直到满足用户的需求。 根据业务逻辑设计自动化系统，驱动正常运行的产品和系统，是一个人类认知上的非凡壮举。

幸运的是，对日益壮大的机器学习科学家群体来说，实现很多任务的自动化并不再屈从于人类所能考虑到的逻辑。 想象一下，假如开发人员要试图解决以下问题之一：

- 编写一个应用程序，接受地理信息、卫星图像和一些历史天气信息，并预测明天的天气；
- 编写一个应用程序，接受自然文本表示的问题，并正确回答该问题；
- 编写一个应用程序，接受一张图像，识别出该图像所包含的人，并在每个人周围绘制轮廓；
- 编写一个应用程序，向用户推荐他们可能喜欢，但在自然浏览过程中不太可能遇到的产品。

在这些情况下，即使是顶级程序员也无法提出完美的解决方案， 原因可能各不相同。有时任务可能遵循一种随着时间推移而变化的模式，我们需要程序来自动调整。 有时任务内的关系可能太复杂（比如像素和抽象类别之间的关系），需要数千或数百万次的计算。 即使人类的眼睛能毫不费力地完成这些难以提出完美解决方案的任务，这其中的计算也超出了人类意识理解范畴。 机器学习（machine learning，ML）是一类强大的可以从经验中学习的技术。 通常采用观测数据或与环境交互的形式，机器学习算法会积累更多的经验，其性能也会逐步提高。 相反，对于刚刚所说的电子商务平台，如果它一直执行相同的业务逻辑，无论积累多少经验，都不会自动提高，除非开发人员认识到问题并更新软件。 本书将带读者开启机器学习之旅，并特别关注深度学习（deep learning，DL）的基础知识。 深度学习是一套强大的技术，它可以推动计算机视觉、自然语言处理、医疗保健和基因组学等不同领域的创新。

 1.1. 日常生活中的机器学习

机器学习应用在日常生活中的方方面面。 现在，假设本书的作者们一起驱车去咖啡店。 阿斯顿拿起一部iPhone，对它说道：“Hey Siri！”手机的语音识别系统就被唤醒了。 接着，李沐对Siri说道：“去星巴克咖啡店。”语音识别系统就自动触发语音转文字功能，并启动地图应用程序， 地图应用程序在启动后筛选了若干条路线，每条路线都显示了预计的通行时间…… 由此可见，机器学习渗透在生活中的方方面面，在短短几秒钟的时间里，人们与智能手机的日常互动就可以涉及几种机器学习模型。

现在，假如需要我们编写程序来响应一个“唤醒词”（比如“Alexa”“小爱同学”和“Hey Siri”）。 我们试着用一台计算机和一个代码编辑器编写代码，如 [图1.1.1]
(https://zh-v2.d2l.ai/chapter\_introduction/index.htmlfig-wake-word)中所示。 问题看似很难解决：麦克风每秒钟将收集大约44000个样本，每个样本都是声波振幅的测量值。而该测量值与唤醒词难以直接关联。那又该如何编写程序，令其输入麦克风采集到的原始音频片段,输出是否{是,否}（表示该片段是否包含唤醒词）的可靠预测呢？我们对编写这个程序毫无头绪，这就是需要机器学习的原因。

%![../\_images/wake-word.svg](https://zh-v2.d2l.ai/\_images/wake-word.svg)

图1.1.1 识别唤醒词

通常，即使我们不知道怎样明确地告诉计算机如何从输入映射到输出，大脑仍然能够自己执行认知功能。 换句话说，即使我们不知道如何编写计算机程序来识别“Alexa”这个词，大脑自己也能够识别它。 有了这一能力，我们就可以收集一个包含大量音频样本的数据集（dataset），并对包含和不包含唤醒词的样本进行标记。 利用机器学习算法，我们不需要设计一个“明确地”识别唤醒词的系统。 相反，我们只需要定义一个灵活的程序算法，其输出由许多参数（parameter）决定，然后使用数据集来确定当下的“最佳参数集”，这些参数通过某种性能度量方式来达到完成任务的最佳性能。

那么到底什么是参数呢？ 参数可以被看作旋钮，旋钮的转动可以调整程序的行为。 任一调整参数后的程序被称为模型（model）。 通过操作参数而生成的所有不同程序（输入-输出映射）的集合称为“模型族”。 使用数据集来选择参数的元程序被称为学习算法（learning algorithm）。

在开始用机器学习算法解决问题之前，我们必须精确地定义问题，确定输入（input）和输出（output）的性质，并选择合适的模型族。 在本例中，模型接收一段音频作为输入，然后在是或否中生成一个选择作为输出。 如果一切顺利，经过一番训练，模型对于“片段是否包含唤醒词”的预测通常是正确的。

现在模型每次听到“Alexa”这个词时都会发出“是”的声音。 由于这里的唤醒词是任意选择的自然语言，因此我们可能需要一个足够丰富的模型族，使模型多元化。 比如，模型族的另一个模型只在听到“Hey Siri”这个词时发出“是”。 理想情况下，同一个模型族应该适合于“Alexa”识别和“Hey Siri”识别，因为从直觉上看，它们似乎是相似的任务。 然而，如果我们想处理完全不同的输入或输出，比如：从图像映射到字幕，或从英语映射到中文，可能需要一个完全不同的模型族。

但如果模型所有的按钮（模型参数）都被随机设置，就不太可能识别出“Alexa”“Hey Siri”或任何其他单词。 在机器学习中，学习（learning）是一个训练模型的过程。 通过这个过程，我们可以发现正确的参数集，从而使模型强制执行所需的行为。 换句话说，我们用数据训练（train）模型。 如 [图1.1.2]
(https://zh-v2.d2l.ai/chapter\_introduction/index.htmlfig-ml-loop)
所示，训练过程通常包含如下步骤：

1. 从一个随机初始化参数的模型开始，这个模型基本没有“智能”；
2. 获取一些数据样本（例如，音频片段以及对应的是或否标签）；
3. 调整参数，使模型在这些样本中表现得更好；
4. 重复第（2）步和第（3）步，直到模型在任务中的表现令人满意。

%![../\_images/ml-loop.svg](https://zh-v2.d2l.ai/\_images/ml-loop.svg)

图1.1.2 一个典型的训练过程

总而言之，我们没有编写唤醒词识别器，而是编写了一个“学习”程序。 如果我们用一个巨大的带标签的数据集，它很可能可以“学习”识别唤醒词。 这种“通过用数据集来确定程序行为”的方法可以被看作用数据编程（programming with data）。 比如，我们可以通过向机器学习系统，提供许多猫和狗的图片来设计一个“猫图检测器”。 检测器最终可以学会：如果输入是猫的图片就输出一个非常大的正数，如果输入是狗的图片就会输出一个非常小的负数。 如果检测器不确定输入的图片中是猫还是狗，它会输出接近于零的数…… 这个例子仅仅是机器学习常见应用的冰山一角， 而深度学习是机器学习的一个主要分支，本节稍后的内容将对其进行更详细的解析。

 1.2. 机器学习中的关键组件

首先介绍一些核心组件。无论什么类型的机器学习问题，都会遇到这些组件：

1. 可以用来学习的数据（data）；
2. 如何转换数据的模型（model）；
3. 一个目标函数（objective function），用来量化模型的有效性；
4. 调整模型参数以优化目标函数的算法（algorithm）。

 1.2.1. 数据

毋庸置疑，如果没有数据，那么数据科学毫无用武之地。 每个数据集由一个个样本（example, sample）组成，大多时候，它们遵循独立同分布(independently and identically distributed, i.i.d.)。 样本有时也叫做数据点（data point）或者数据实例（data instance），通常每个样本由一组称为特征（features，或协变量（covariates））的属性组成。 机器学习模型会根据这些属性进行预测。 在上面的监督学习问题中，要预测的是一个特殊的属性，它被称为标签（label，或目标（target））。

当处理图像数据时，每一张单独的照片即为一个样本，它的特征由每个像素数值的有序列表表示。 比如，200×200彩色照片由200×200×3=120000个数值组成，其中的“3”对应于每个空间位置的红、绿、蓝通道的强度。 再比如，对于一组医疗数据，给定一组标准的特征（如年龄、生命体征和诊断），此数据可以用来尝试预测患者是否会存活。

当每个样本的特征类别数量都是相同的时候，其特征向量是固定长度的，这个长度被称为数据的维数（dimensionality）。 固定长度的特征向量是一个方便的属性，它可以用来量化学习大量样本。

然而，并不是所有的数据都可以用“固定长度”的向量表示。 以图像数据为例，如果它们全部来自标准显微镜设备，那么“固定长度”是可取的； 但是如果图像数据来自互联网，它们很难具有相同的分辨率或形状。 这时，将图像裁剪成标准尺寸是一种方法，但这种办法很局限，有丢失信息的风险。 此外，文本数据更不符合“固定长度”的要求。 比如，对于亚马逊等电子商务网站上的客户评论，有些文本数据很简短（比如“好极了”），有些则长篇大论。 与传统机器学习方法相比，深度学习的一个主要优势是可以处理不同长度的数据。

一般来说，拥有越多数据的时候，工作就越容易。 更多的数据可以被用来训练出更强大的模型，从而减少对预先设想假设的依赖。 数据集的由小变大为现代深度学习的成功奠定基础。 在没有大数据集的情况下，许多令人兴奋的深度学习模型黯然失色。 就算一些深度学习模型在小数据集上能够工作，但其效能并不比传统方法高。

请注意，仅仅拥有海量的数据是不够的，我们还需要正确的数据。 如果数据中充满了错误，或者如果数据的特征不能预测任务目标，那么模型很可能无效。 有一句古语很好地反映了这个现象：“输入的是垃圾，输出的也是垃圾。”（“Garbage in, garbage out.”） 此外，糟糕的预测性能甚至会加倍放大事态的严重性。 在一些敏感应用中，如预测性监管、简历筛选和用于贷款的风险模型，我们必须特别警惕垃圾数据带来的后果。 一种常见的问题来自不均衡的数据集，比如在一个有关医疗的训练数据集中，某些人群没有样本表示。 想象一下，假设我们想要训练一个皮肤癌识别模型，但它（在训练数据集中）从未“见过”黑色皮肤的人群，这个模型就会顿时束手无策。

再比如，如果用“过去的招聘决策数据”来训练一个筛选简历的模型，那么机器学习模型可能会无意中捕捉到历史残留的不公正，并将其自动化。 然而，这一切都可能在不知情的情况下发生。 因此，当数据不具有充分代表性，甚至包含了一些社会偏见时，模型就很有可能有偏见。

 1.2.2. 模型

大多数机器学习会涉及到数据的转换。 比如一个“摄取照片并预测笑脸”的系统。再比如通过摄取到的一组传感器读数预测读数的正常与异常程度。 虽然简单的模型能够解决如上简单的问题，但本书中关注的问题超出了经典方法的极限。 深度学习与经典方法的区别主要在于：前者关注的功能强大的模型，这些模型由神经网络错综复杂的交织在一起，包含层层数据转换，因此被称为深度学习（deep learning）。 在讨论深度模型的过程中，本书也将提及一些传统方法。

 1.2.3. 目标函数

前面的内容将机器学习介绍为“从经验中学习”。 这里所说的“学习”，是指自主提高模型完成某些任务的效能。 但是，什么才算真正的提高呢？ 在机器学习中，我们需要定义模型的优劣程度的度量，这个度量在大多数情况是“可优化”的，这被称之为目标函数（objective function）。 我们通常定义一个目标函数，并希望优化它到最低点。 因为越低越好，所以这些函数有时被称为损失函数（loss function，或cost function）。 但这只是一个惯例，我们也可以取一个新的函数，优化到它的最高点。 这两个函数本质上是相同的，只是翻转一下符号。

当任务在试图预测数值时，最常见的损失函数是平方误差（squared error），即预测值与实际值之差的平方。 当试图解决分类问题时，最常见的目标函数是最小化错误率，即预测与实际情况不符的样本比例。 有些目标函数（如平方误差）很容易被优化，有些目标（如错误率）由于不可微性或其他复杂性难以直接优化。 在这些情况下，通常会优化替代目标。

通常，损失函数是根据模型参数定义的，并取决于数据集。 在一个数据集上，我们可以通过最小化总损失来学习模型参数的最佳值。 该数据集由一些为训练而收集的样本组成，称为训练数据集（training dataset，或称为训练集（training set））。 然而，在训练数据上表现良好的模型，并不一定在“新数据集”上有同样的性能，这里的“新数据集”通常称为测试数据集（test dataset，或称为测试集（test set））。

综上所述，可用数据集通常可以分成两部分：训练数据集用于拟合模型参数，测试数据集用于评估拟合的模型。 然后我们观察模型在这两部分数据集的性能。 “一个模型在训练数据集上的性能”可以被想象成“一个学生在模拟考试中的分数”。 这个分数用来为一些真正的期末考试做参考，即使成绩令人鼓舞，也不能保证期末考试成功。 换言之，测试性能可能会显著偏离训练性能。 当一个模型在训练集上表现良好，但不能推广到测试集时，这个模型被称为过拟合（overfitting）的。 就像在现实生活中，尽管模拟考试考得很好，真正的考试不一定百发百中。

 1.2.4. 优化算法

当我们获得了一些数据源及其表示、一个模型和一个合适的损失函数，接下来就需要一种算法，它能够搜索出最佳参数，以最小化损失函数。 深度学习中，大多流行的优化算法通常基于一种基本方法–梯度下降（gradient descent）。 简而言之，在每个步骤中，梯度下降法都会检查每个参数，看看如果仅对该参数进行少量变动，训练集损失会朝哪个方向移动。 然后，它在可以减少损失的方向上优化参数。

 1.3. 各种机器学习问题

在机器学习的广泛应用中，唤醒词问题只是冰山一角。 前面唤醒词识别的例子，只是机器学习可以解决的众多问题中的一个。 下面将列出一些常见的机器学习问题和应用，为之后本书的讨论做铺垫。 接下来会经常引用前面提到的概念，如数据、模型和优化算法。

 1.3.1. 监督学习

监督学习（supervised learning）擅长在“给定输入特征”的情况下预测标签。 每个“特征-标签”对都称为一个样本（example）。 有时，即使标签是未知的，样本也可以指代输入特征。 我们的目标是生成一个模型，能够将任何输入特征映射到标签（即预测）。

举一个具体的例子： 假设我们需要预测患者的心脏病是否会发作，那么观察结果“心脏病发作”或“心脏病没有发作”将是样本的标签。 输入特征可能是生命体征，如心率、舒张压和收缩压等。

监督学习之所以能发挥作用，是因为在训练参数时，我们为模型提供了一个数据集，其中每个样本都有真实的标签。 用概率论术语来说，我们希望预测“估计给定输入特征的标签”的条件概率。 虽然监督学习只是几大类机器学习问题之一，但是在工业中，大部分机器学习的成功应用都使用了监督学习。 这是因为在一定程度上，许多重要的任务可以清晰地描述为，在给定一组特定的可用数据的情况下，估计未知事物的概率。比如：

- 根据计算机断层扫描（Computed Tomography，CT）肿瘤图像，预测是否为癌症；
- 给出一个英语句子，预测正确的法语翻译；
- 根据本月的财务报告数据，预测下个月股票的价格；

监督学习的学习过程一般可以分为三大步骤：

1. 从已知大量数据样本中随机选取一个子集，为每个样本获取真实标签。有时，这些样本已有标签（例如，患者是否在下一年内康复？）；有时，这些样本可能需要被人工标记（例如，图像分类）。这些输入和相应的标签一起构成了训练数据集；
2. 选择有监督的学习算法，它将训练数据集作为输入，并输出一个“已完成学习的模型”；
3. 将之前没有见过的样本特征放到这个“已完成学习的模型”中，使用模型的输出作为相应标签的预测。

整个监督学习过程如 [图1.3.1]
(https://zh-v2.d2l.ai/chapter\_introduction/index.htmlfig-supervised-learning) 
所示。

![../\_images/supervised-learning.svg](https://zh-v2.d2l.ai/\_images/supervised-learning.svg)

图1.3.1 监督学习

综上所述，即使使用简单的描述给定输入特征的预测标签，监督学习也可以采取多种形式的模型，并且需要大量不同的建模决策，这取决于输入和输出的类型、大小和数量。 例如，我们使用不同的模型来处理“任意长度的序列”或“固定长度的序列”。

 1.3.1.1. 回归

回归（regression）是最简单的监督学习任务之一。 假设有一组房屋销售数据表格，其中每行对应一个房子，每列对应一个相关的属性，例如房屋的面积、卧室的数量、浴室的数量以及到镇中心的步行距离，等等。 每一行的属性构成了一个房子样本的特征向量。 如果一个人住在纽约或旧金山，而且他不是亚马逊、谷歌、微软或Facebook的首席执行官，那么他家的特征向量（房屋面积，卧室数量，浴室数量，步行距离）可能类似于：[600,1,1,60]。 如果一个人住在匹兹堡，这个特征向量可能更接近[3000,4,3,10]…… 当人们在市场上寻找新房子时，可能需要估计一栋房子的公平市场价值。 为什么这个任务可以归类为回归问题呢？本质上是输出决定的。 销售价格（即标签）是一个数值。 当标签取任意数值时，我们称之为回归问题，此时的目标是生成一个模型，使它的预测非常接近实际标签值。

生活中的许多问题都可归类为回归问题。 比如，预测用户对一部电影的评分可以被归类为一个回归问题。 这里有一个小插曲：在2009年，如果有人设计了一个很棒的算法来预测电影评分，那可能会赢得[100万美元的奈飞奖](https://en.wikipedia.org/wiki/Netflix\_Prize)。 再比如，预测病人在医院的住院时间也是一个回归问题。 总而言之，判断回归问题的一个很好的经验法则是，任何有关“有多少”的问题很可能就是回归问题。比如：

- 这个手术需要多少小时；
- 在未来6小时，这个镇会有多少降雨量。

即使你以前从未使用过机器学习，可能在不经意间，已经解决了一些回归问题。 例如，你让人修理了排水管，承包商花了3小时清除污水管道中的污物，然后他寄给你一张350美元的账单。 而你的朋友雇了同一个承包商2小时，他收到了250美元的账单。 如果有人请你估算清理污物的费用，你可以假设承包商收取一些基本费用，然后按小时收费。 如果这些假设成立，那么给出这两个数据样本，你就已经可以确定承包商的定价结构：50美元上门服务费，另外每小时100美元。 在不经意间，你就已经理解并应用了线性回归算法。

然而，以上假设有时并不可取。 例如，一些差异是由于两个特征之外的几个因素造成的。 在这些情况下，我们将尝试学习最小化“预测值和实际标签值的差异”的模型。 本书大部分章节将关注平方误差损失函数的最小化。

 1.3.1.2. 分类

虽然回归模型可以很好地解决“有多少”的问题，但是很多问题并非如此。 例如，一家银行希望在其移动应用程序中添加支票扫描功能。 具体地说，这款应用程序能够自动理解从图像中看到的文本，并将手写字符映射到对应的已知字符之上。 这种“哪一个”的问题叫做分类（classification）问题。 分类问题希望模型能够预测样本属于哪个类别（category，正式称为类（class））。 例如，手写数字可能有10类，标签被设置为数字0～9。 最简单的分类问题是只有两类，这被称之为二项分类（binomial classification）。 例如，数据集可能由动物图像组成，标签可能是猫狗{猫,狗}两类。 回归是训练一个回归函数来输出一个数值； 分类是训练一个分类器来输出预测的类别。

然而模型怎么判断得出这种“是”或“不是”的硬分类预测呢？ 我们可以试着用概率语言来理解模型。 给定一个样本特征，模型为每个可能的类分配一个概率。 比如，之前的猫狗分类例子中，分类器可能会输出图像是猫的概率为0.9。 0.9这个数字表达什么意思呢？ 可以这样理解：分类器90%确定图像描绘的是一只猫。 预测类别的概率的大小传达了一种模型的不确定性，本书后面章节将讨论其他运用不确定性概念的算法。

当有两个以上的类别时，我们把这个问题称为多项分类（multiclass classification）问题。 常见的例子包括手写字符识别 {0,1,2,...9,a,b,c,...}。 与解决回归问题不同，分类问题的常见损失函数被称为交叉熵（cross-entropy），本书 [3.4节]
(https://zh-v2.d2l.ai/chapter\_linear-networks/softmax-regression.htmlsec-softmax) 
将详细阐述。

请注意，最常见的类别不一定是最终用于决策的类别。 举个例子，假设后院有一个如 [图1.3.2]
(https://zh-v2.d2l.ai/chapter\_introduction/index.htmlfig-death-cap) 
所示的蘑菇。


%[![../\_images/death-cap.jpg](https://zh-v2.d2l.ai/\_images/death-cap.jpg)](https://zh-v2.d2l.ai/\_images/death-cap.jpg)
% 图1.3.2 死帽蕈——不能吃！！

现在，我们想要训练一个毒蘑菇检测分类器，根据照片预测蘑菇是否有毒。 假设这个分类器输出 [图1.3.2]
(https://zh-v2.d2l.ai/chapter\_introduction/index.htmlfig-death-cap) 
包含死帽蕈的概率是0.2。 换句话说，分类器80\%确定图中的蘑菇不是死帽蕈。 尽管如此，我们也不会吃它，因为不值得冒20\%的死亡风险。 换句话说，不确定风险的影响远远大于收益。 因此，我们需要将“预期风险”作为损失函数，即需要将结果的概率乘以与之相关的收益（或伤害）。 在这种情况下，食用蘑菇造成的损失为0.2×∞+0.8×0=∞，而丢弃蘑菇的损失为0.2×0+0.8×1=0.8。 事实上，谨慎是有道理的， [图1.3.2]
(https://zh-v2.d2l.ai/chapter\_introduction/index.htmlfig-death-cap)
中的蘑菇实际上是一个死帽蕈。

分类可能变得比二项分类、多项分类复杂得多。 例如，有一些分类任务的变体可以用于寻找层次结构，层次结构假定在许多类之间存在某种关系。 因此，并不是所有的错误都是均等的。 人们宁愿错误地分入一个相关的类别，也不愿错误地分入一个遥远的类别，这通常被称为层次分类(hierarchical classification)。 早期的一个例子是[卡尔·林奈]
(https://en.wikipedia.org/wiki/Carl\_Linnaeus)，
他对动物进行了层次分类。

在动物分类的应用中，把一只狮子狗误认为雪纳瑞可能不会太糟糕。 但如果模型将狮子狗与恐龙混淆，就滑稽至极了。 层次结构相关性可能取决于模型的使用者计划如何使用模型。 例如，响尾蛇和乌梢蛇血缘上可能很接近，但如果把响尾蛇误认为是乌梢蛇可能会是致命的。 因为响尾蛇是有毒的，而乌梢蛇是无毒的。

 1.3.1.3. 标记问题

有些分类问题很适合于二项分类或多项分类。 例如，我们可以训练一个普通的二项分类器来区分猫和狗。 运用最前沿的计算机视觉的算法，这个模型可以很轻松地被训练。 尽管如此，无论模型有多精确，当分类器遇到新的动物时可能会束手无策。 比如 [图1.3.3](https://zh-v2.d2l.ai/chapter\_introduction/index.htmlfig-stackedanimals)所示的这张“不来梅的城市音乐家”的图像 （这是一个流行的德国童话故事），图中有一只猫、一只公鸡、一只狗、一头驴，背景是一些树。 取决于我们最终想用模型做什么，将其视为二项分类问题可能没有多大意义。 取而代之，我们可能想让模型描绘输入图像的内容，一只猫、一只公鸡、一只狗，还有一头驴。

[![../\_images/stackedanimals.png](https://zh-v2.d2l.ai/\_images/stackedanimals.png)](https://zh-v2.d2l.ai/\_images/stackedanimals.png)

图1.3.3 一只猫、一只公鸡、一只狗、一头驴

学习预测不相互排斥的类别的问题称为多标签分类（multi-label classification）。 举个例子，人们在技术博客上贴的标签，比如“机器学习”“技术”“小工具”“编程语言”“Linux”“云计算”“AWS”。 一篇典型的文章可能会用5～10个标签，因为这些概念是相互关联的。 关于“云计算”的帖子可能会提到“AWS”，而关于“机器学习”的帖子也可能涉及“编程语言”。

此外，在处理生物医学文献时，我们也会遇到这类问题。 正确地标记文献很重要，有利于研究人员对文献进行详尽的审查。 在美国国家医学图书馆（The United States National Library of Medicine），一些专业的注释员会检查每一篇在PubMed中被索引的文章，以便将其与Mesh中的相关术语相关联（Mesh是一个大约有28000个标签的集合）。 这是一个十分耗时的过程，注释器通常在归档和标记之间有一年的延迟。 这里，机器学习算法可以提供临时标签，直到每一篇文章都有严格的人工审核。 事实上，近几年来，BioASQ组织已经[举办比赛](http://bioasq.org/)来完成这项工作。

 1.3.1.4. 搜索

有时，我们不仅仅希望输出一个类别或一个实值。 在信息检索领域，我们希望对一组项目进行排序。 以网络搜索为例，目标不是简单的“查询（query）-网页（page）”分类，而是在海量搜索结果中找到用户最需要的那部分。 搜索结果的排序也十分重要，学习算法需要输出有序的元素子集。 换句话说，如果要求我们输出字母表中的前5个字母，返回“A、B、C、D、E”和“C、A、B、E、D”是不同的。 即使结果集是相同的，集内的顺序有时却很重要。

该问题的一种可能的解决方案：首先为集合中的每个元素分配相应的相关性分数，然后检索评级最高的元素。[PageRank](https://en.wikipedia.org/wiki/PageRank)，谷歌搜索引擎背后最初的秘密武器就是这种评分系统的早期例子，但它的奇特之处在于它不依赖于实际的查询。 在这里，他们依靠一个简单的相关性过滤来识别一组相关条目，然后根据PageRank对包含查询条件的结果进行排序。 如今，搜索引擎使用机器学习和用户行为模型来获取网页相关性得分，很多学术会议也致力于这一主题。

 1.3.1.5. 推荐系统

另一类与搜索和排名相关的问题是推荐系统（recommender system），它的目标是向特定用户进行“个性化”推荐。 例如，对于电影推荐，科幻迷和喜剧爱好者的推荐结果页面可能会有很大不同。 类似的应用也会出现在零售产品、音乐和新闻推荐等等。

在某些应用中，客户会提供明确反馈，表达他们对特定产品的喜爱程度。 例如，亚马逊上的产品评级和评论。 在其他一些情况下，客户会提供隐性反馈。 例如，某用户跳过播放列表中的某些歌曲，这可能说明这些歌曲对此用户不大合适。 总的来说，推荐系统会为“给定用户和物品”的匹配性打分，这个“分数”可能是估计的评级或购买的概率。 由此，对于任何给定的用户，推荐系统都可以检索得分最高的对象集，然后将其推荐给用户。以上只是简单的算法，而工业生产的推荐系统要先进得多，它会将详细的用户活动和项目特征考虑在内。 推荐系统算法经过调整，可以捕捉一个人的偏好。 比如， [图1.3.4]
(https://zh-v2.d2l.ai/chapter\_introduction/index.htmlfig-deeplearning-amazon) 
是亚马逊基于个性化算法推荐的深度学习书籍，成功地捕捉了作者的喜好。

%![../\_images/deeplearning-amazon.jpg](https://zh-v2.d2l.ai/\_images/deeplearning-amazon.jpg)

图1.3.4 亚马逊推荐的深度学习书籍

尽管推荐系统具有巨大的应用价值，但单纯用它作为预测模型仍存在一些缺陷。 首先，我们的数据只包含“审查后的反馈”：用户更倾向于给他们感觉强烈的事物打分。 例如，在五分制电影评分中，会有许多五星级和一星级评分，但三星级却明显很少。 此外，推荐系统有可能形成反馈循环：推荐系统首先会优先推送一个购买量较大（可能被认为更好）的商品，然而目前用户的购买习惯往往是遵循推荐算法，但学习算法并不总是考虑到这一细节，进而更频繁地被推荐。 综上所述，关于如何处理审查、激励和反馈循环的许多问题，都是重要的开放性研究问题。

 1.3.1.6. 序列学习

以上大多数问题都具有固定大小的输入和产生固定大小的输出。 例如，在预测房价的问题中，我们考虑从一组固定的特征：房屋面积、卧室数量、浴室数量、步行到市中心的时间； 图像分类问题中，输入为固定尺寸的图像，输出则为固定数量（有关每一个类别）的预测概率； 在这些情况下，模型只会将输入作为生成输出的“原料”，而不会“记住”输入的具体内容。

如果输入的样本之间没有任何关系，以上模型可能完美无缺。 但是如果输入是连续的，模型可能就需要拥有“记忆”功能。 比如，我们该如何处理视频片段呢？ 在这种情况下，每个视频片段可能由不同数量的帧组成。 通过前一帧的图像，我们可能对后一帧中发生的事情更有把握。 语言也是如此，机器翻译的输入和输出都为文字序列。

再比如，在医学上序列输入和输出就更为重要。 设想一下，假设一个模型被用来监控重症监护病人，如果他们在未来24小时内死亡的风险超过某个阈值，这个模型就会发出警报。 我们绝不希望抛弃过去每小时有关病人病史的所有信息，而仅根据最近的测量结果做出预测。

这些问题是序列学习的实例，是机器学习最令人兴奋的应用之一。 序列学习需要摄取输入序列或预测输出序列，或两者兼而有之。 具体来说，输入和输出都是可变长度的序列，例如机器翻译和从语音中转录文本。 虽然不可能考虑所有类型的序列转换，但以下特殊情况值得一提。

标记和解析。
这涉及到用属性注释文本序列。 换句话说，输入和输出的数量基本上是相同的。 例如，我们可能想知道动词和主语在哪里，或者可能想知道哪些单词是命名实体。 通常，目标是基于结构和语法假设对文本进行分解和注释，以获得一些注释。 这听起来比实际情况要复杂得多。 下面是一个非常简单的示例，它使用“标记”来注释一个句子，该标记指示哪些单词引用命名实体。 标记为“Ent”，是实体（entity）的简写。

```
Tom has dinner in Washington with Sally
Ent  -    -    -     Ent      -    Ent
```

自动语音识别。在语音识别中，输入序列是说话人的录音（如 [图1.3.5]
(https://zh-v2.d2l.ai/chapter\_introduction/index.htmlfig-speech) 
所示），输出序列是说话人所说内容的文本记录。 它的挑战在于，与文本相比，音频帧多得多（声音通常以8kHz或16kHz采样）。 也就是说，音频和文本之间没有1:1的对应关系，因为数千个样本可能对应于一个单独的单词。 这也是“序列到序列”的学习问题，其中输出比输入短得多。

[![../\_images/speech.png](https://zh-v2.d2l.ai/\_images/speech.png)](https://zh-v2.d2l.ai/\_images/speech.png)

图1.3.5 `-D-e-e-p- L-ea-r-ni-ng-` 在录音中。

文本到语音。这与自动语音识别相反。 换句话说，输入是文本，输出是音频文件。 在这种情况下，输出比输入长得多。 虽然人类很容易识判断发音别扭的音频文件，但这对计算机来说并不是那么简单。

机器翻译。 在语音识别中，输入和输出的出现顺序基本相同。 而在机器翻译中，颠倒输入和输出的顺序非常重要。 换句话说，虽然我们仍将一个序列转换成另一个序列，但是输入和输出的数量以及相应序列的顺序大都不会相同。 比如下面这个例子，“错误的对齐”反应了德国人喜欢把动词放在句尾的特殊倾向。

```
德语:           Haben Sie sich schon dieses grossartige Lehrwerk angeschaut?
英语:          Did you already check out this excellent tutorial?
错误的对齐:  Did you yourself already this excellent tutorial looked-at?
```

其他学习任务也有序列学习的应用。 例如，确定“用户阅读网页的顺序”是二维布局分析问题。 再比如，对话问题对序列的学习更为复杂：确定下一轮对话，需要考虑对话历史状态以及现实世界的知识…… 如上这些都是热门的序列学习研究领域。

 1.3.2. 无监督学习

到目前为止，所有的例子都与监督学习有关，即需要向模型提供巨大数据集：每个样本包含特征和相应标签值。 打趣一下，“监督学习”模型像一个打工仔，有一份极其专业的工作和一位极其平庸的老板。 老板站在身后，准确地告诉模型在每种情况下应该做什么，直到模型学会从情况到行动的映射。 取悦这位老板很容易，只需尽快识别出模式并模仿他们的行为即可。

相反，如果工作没有十分具体的目标，就需要“自发”地去学习了。 比如，老板可能会给我们一大堆数据，然后要求用它做一些数据科学研究，却没有对结果有要求。 这类数据中不含有“目标”的机器学习问题通常被为无监督学习（unsupervised learning）， 本书后面的章节将讨论无监督学习技术。 那么无监督学习可以回答什么样的问题呢？来看看下面的例子。

- 聚类（clustering）问题：没有标签的情况下，我们是否能给数据分类呢？比如，给定一组照片，我们能把它们分成风景照片、狗、婴儿、猫和山峰的照片吗？同样，给定一组用户的网页浏览记录，我们能否将具有相似行为的用户聚类呢？
- 主成分分析（principal component analysis）问题：我们能否找到少量的参数来准确地捕捉数据的线性相关属性？比如，一个球的运动轨迹可以用球的速度、直径和质量来描述。再比如，裁缝们已经开发出了一小部分参数，这些参数相当准确地描述了人体的形状，以适应衣服的需要。另一个例子：在欧几里得空间中是否存在一种（任意结构的）对象的表示，使其符号属性能够很好地匹配?这可以用来描述实体及其关系，例如“罗马” − “意大利” + “法国” = “巴黎”。
- 因果关系（causality）和概率图模型（probabilistic graphical models）问题：我们能否描述观察到的许多数据的根本原因？例如，如果我们有关于房价、污染、犯罪、地理位置、教育和工资的人口统计数据，我们能否简单地根据经验数据发现它们之间的关系？
- 生成对抗性网络（generative adversarial networks）：为我们提供一种合成数据的方法，甚至像图像和音频这样复杂的非结构化数据。潜在的统计机制是检查真实和虚假数据是否相同的测试，它是无监督学习的另一个重要而令人兴奋的领域。

 1.3.3. 与环境互动

有人一直心存疑虑：机器学习的输入（数据）来自哪里？机器学习的输出又将去往何方？ 到目前为止，不管是监督学习还是无监督学习，我们都会预先获取大量数据，然后启动模型，不再与环境交互。 这里所有学习都是在算法与环境断开后进行的，被称为离线学习（offline learning）。 对于监督学习，从环境中收集数据的过程类似于 [图1.3.6]
(https://zh-v2.d2l.ai/chapter\_introduction/index.htmlfig-data-collection)。

%![../\_images/data-collection.svg](https://zh-v2.d2l.ai/\_images/data-collection.svg)

图1.3.6 从环境中为监督学习收集数据。

这种简单的离线学习有它的魅力。 好的一面是，我们可以孤立地进行模式识别，而不必分心于其他问题。 但缺点是，解决的问题相当有限。 这时我们可能会期望人工智能不仅能够做出预测，而且能够与真实环境互动。 与预测不同，“与真实环境互动”实际上会影响环境。 这里的人工智能是“智能代理”，而不仅是“预测模型”。 因此，我们必须考虑到它的行为可能会影响未来\\观察结果。

考虑“与真实环境互动”将打开一整套新的建模问题。以下只是几个例子。

- 环境还记得我们以前做过什么吗？
- 环境是否有助于我们建模？例如，用户将文本读入语音识别器。
- 环境是否想要打败模型？例如，一个对抗性的设置，如垃圾邮件过滤或玩游戏？
- 环境是否重要？
- 环境是否变化？例如，未来的数据是否总是与过去相似，还是随着时间的推移会发生变化？是自然变化还是响应我们的自动化工具而发生变化？

当训练和测试数据不同时，最后一个问题提出了分布偏移（distribution shift）的问题。 接下来的内容将简要描述强化学习问题，这是一类明确考虑与环境交互的问题。

 1.3.4. 强化学习

如果你对使用机器学习开发与环境交互并采取行动感兴趣，那么最终可能会专注于强化学习（reinforcement learning）。 这可能包括应用到机器人、对话系统，甚至开发视频游戏的人工智能（AI）。 深度强化学习（deep reinforcement learning）将深度学习应用于强化学习的问题，是非常热门的研究领域。 突破性的深度Q网络（Q-network）在雅达利游戏中仅使用视觉输入就击败了人类， 以及 AlphaGo 程序在棋盘游戏围棋中击败了世界冠军，是两个突出强化学习的例子。

在强化学习问题中，智能体（agent）在一系列的时间步骤上与环境交互。 在每个特定时间点，智能体从环境接收一些观察（observation），并且必须选择一个动作（action），然后通过某种机制（有时称为执行器）将其传输回环境，最后智能体从环境中获得奖励（reward）。 此后新一轮循环开始，智能体接收后续观察，并选择后续操作，依此类推。 强化学习的过程在 [图1.3.7]
(https://zh-v2.d2l.ai/chapter\_introduction/index.htmlfig-rl-environment) 
中进行了说明。 请注意，强化学习的目标是产生一个好的策略（policy）。 强化学习智能体选择的“动作”受策略控制，即一个从环境观察映射到行动的功能。

%![../\_images/rl-environment.svg](https://zh-v2.d2l.ai/\_images/rl-environment.svg)

图1.3.7 强化学习和环境之间的相互作用

强化学习框架的通用性十分强大。 例如，我们可以将任何监督学习问题转化为强化学习问题。 假设我们有一个分类问题，可以创建一个强化学习智能体，每个分类对应一个“动作”。 然后，我们可以创建一个环境，该环境给予智能体的奖励。 这个奖励与原始监督学习问题的损失函数是一致的。

当然，强化学习还可以解决许多监督学习无法解决的问题。 例如，在监督学习中，我们总是希望输入与正确的标签相关联。 但在强化学习中，我们并不假设环境告诉智能体每个观测的最优动作。 一般来说，智能体只是得到一些奖励。 此外，环境甚至可能不会告诉是哪些行为导致了奖励。

以强化学习在国际象棋的应用为例。 唯一真正的奖励信号出现在游戏结束时：当智能体获胜时，智能体可以得到奖励1；当智能体失败时，智能体将得到奖励-1。 因此，强化学习者必须处理学分分配（credit assignment）问题：决定哪些行为是值得奖励的，哪些行为是需要惩罚的。 就像一个员工升职一样，这次升职很可能反映了前一年的大量的行动。 要想在未来获得更多的晋升，就需要弄清楚这一过程中哪些行为导致了晋升。

强化学习可能还必须处理部分可观测性问题。 也就是说，当前的观察结果可能无法阐述有关当前状态的所有信息。 比方说，一个清洁机器人发现自己被困在一个许多相同的壁橱的房子里。 推断机器人的精确位置（从而推断其状态），需要在进入壁橱之前考虑它之前的观察结果。

最后，在任何时间点上，强化学习智能体可能知道一个好的策略，但可能有许多更好的策略从未尝试过的。 强化学习智能体必须不断地做出选择：是应该利用当前最好的策略，还是探索新的策略空间（放弃一些短期回报来换取知识）。

一般的强化学习问题是一个非常普遍的问题。 智能体的动作会影响后续的观察，而奖励只与所选的动作相对应。 环境可以是完整观察到的，也可以是部分观察到的,解释所有这些复杂性可能会对研究人员要求太高。 此外，并不是每个实际问题都表现出所有这些复杂性。 因此，学者们研究了一些特殊情况下的强化学习问题。

当环境可被完全观察到时，强化学习问题被称为马尔可夫决策过程（markov decision process）。 当状态不依赖于之前的操作时，我们称该问题为上下文赌博机（contextual bandit problem）。 当没有状态，只有一组最初未知回报的可用动作时，这个问题就是经典的多臂赌博机（multi-armed bandit problem）。

 1.4. 起源

为了解决各种各样的机器学习问题，深度学习提供了强大的工具。 虽然许多深度学习方法都是最近才有重大突破，但使用数据和神经网络编程的核心思想已经研究了几个世纪。 事实上，人类长期以来就有分析数据和预测未来结果的愿望，而自然科学大部分都植根于此。 例如，伯努利分布是以[雅各布•伯努利（1654-1705）]
(https://en.wikipedia.org/wiki/JacobuBernoulli)
命名的。 而高斯分布是由[卡尔•弗里德里希•高斯（1777-1855）]
(https://en.wikipedia.org/wiki/Carl\_Friedrich\_Gauss)
发现的， 他发明了最小均方算法，至今仍用于解决从保险计算到医疗诊断的许多问题。 这些工具算法催生了自然科学中的一种实验方法——例如，电阻中电流和电压的欧姆定律可以用线性模型完美地描述。

即使在中世纪，数学家对估计（estimation）也有敏锐的直觉。 例如，[雅各布·克贝尔 (1460–1533)]
(https://www.maa.org/press/periodicals/convergence/mathematical-treasures-jacob-kobels-geometry)
的几何学书籍举例说明，通过平均16名成年男性的脚的长度，可以得出一英尺的长度。

[![../\_images/koebel.jpg](https://zh-v2.d2l.ai/\_images/koebel.jpg)](https://zh-v2.d2l.ai/\_images/koebel.jpg)

图1.4.1 估计一英尺的长度

[图1.4.1]
(https://zh-v2.d2l.ai/chapter\_introduction/index.htmlfig-koebel) 
说明了这个估计器是如何工作的。 16名成年男子被要求脚连脚排成一行。 然后将它们的总长度除以16，得到现在等于1英尺的估计值。 

这个算法后来被改进以处理畸形的脚——将拥有最短和最长脚的两个人送走，对其余的人取平均值。 这是最早的修剪均值估计的例子之一。

随着数据的收集和可获得性，统计数据真正实现了腾飞。 [罗纳德·费舍尔（1890-1962）]
%(https://en.wikipedia.org/wiki/Ronald\_-Fisher)
对统计理论和在遗传学中的应用做出了重大贡献。 他的许多算法（如线性判别分析）和公式（如费舍尔信息矩阵）至今仍被频繁使用。 

甚至，费舍尔在1936年发布的鸢尾花卉数据集，有时仍然被用来解读机器学习算法。 

他也是优生学的倡导者，这提醒我们：数据科学在道德上存疑的使用，与其在工业和自然科学中的生产性使用一样，有着悠远而持久的历史。

机器学习的第二个影响来自[克劳德·香农(1916–2001)] %(https://en.wikipedia.org/wiki/Claude\_Shannon)
的信息论和[艾伦·图灵（1912-1954）] %(https://en.wikipedia.org/wiki/Alan\_Turing)
的计算理论。 图灵在他著名的论文《计算机器与智能》 ([Turing, 1950] %(https://zh-v2.d2l.ai/chapter\_references/zreferences.htmlid171)) 
中提出了“机器能思考吗？”的问题。 在他所描述的图灵测试中，如果人类评估者很难根据文本互动区分机器和人类的回答，那么机器就可以被认为是“智能的”。

另一个影响可以在神经科学和心理学中找到。 其中，最古老的算法之一是[唐纳德·赫布 (1904–1985)]
(https://en.wikipedia.org/wiki/Donald\_O.\_Hebb)
开创性的著作《行为的组织》 ([Hebb and Hebb, 1949]
(https://zh-v2.d2l.ai/chapter\_references/zreferences.htmlid62)) 。
他提出神经元通过积极强化学习，是Rosenblatt感知器学习算法的原型，被称为“赫布学习”。 这个算法也为当今深度学习的许多随机梯度下降算法奠定了基础：强化期望行为和减少不良行为，从而在神经网络中获得良好的参数设置。

神经网络（neural networks）的得名源于生物灵感。 一个多世纪以来（追溯到1873年亚历山大·贝恩和1890年詹姆斯·谢林顿的模型），研究人员一直试图组装类似于相互作用的神经元网络的计算电路。 随着时间的推移，对生物学的解释变得不再肤浅，但这个名字仍然存在。 其核心是当今大多数网络中都可以找到的几个关键原则：

- 线性和非线性处理单元的交替，通常称为层（layers）；
- 使用链式规则（也称为反向传播（backpropagation））一次性调整网络中的全部参数。

经过最初的快速发展，神经网络的研究从1995年左右开始停滞不前，直到2005年才稍有起色。 这主要是因为两个原因。 首先，训练网络（在计算上）非常昂贵。 在上个世纪末，随机存取存储器（RAM）非常强大，而计算能力却很弱。 其次，数据集相对较小。 事实上，费舍尔1932年的鸢尾花卉数据集是测试算法有效性的流行工具， 而MNIST数据集的60000个手写数字的数据集被认为是巨大的。 考虑到数据和计算的稀缺性，核方法（kernel method）、决策树（decision tree）和图模型（graph models）等强大的统计工具（在经验上）证明是更为优越的。 与神经网络不同的是，这些算法不需要数周的训练，而且有很强的理论依据，可以提供可预测的结果。

 1.5. 深度学习的发展

大约2010年开始，那些在计算上看起来不可行的神经网络算法变得热门起来，实际上是以下两点导致的： 其一，随着互联网的公司的出现，为数亿在线用户提供服务，大规模数据集变得触手可及； 另外，廉价又高质量的传感器、廉价的数据存储（克莱德定律）以及廉价计算（摩尔定律）的普及，特别是GPU的普及，使大规模算力唾手可得。

这一点在 [表1.5.1](https://zh-v2.d2l.ai/chapter\_introduction/index.htmltab-intro-decade) 
中得到了说明。



| 年代 |              数据规模 |   内存 |        每秒浮点运算 |
| ---: | --------------------: | -----: | ------------------: |
| 1970 |      100 （鸢尾花卉） |   1 KB | 100 KF (Intel 8080) |
| 1980 |    1 K （波士顿房价） | 100 KB |  1 MF (Intel 80186) |
| 1990 | 10 K （光学字符识别） |  10 MB | 10 MF (Intel 80486) |
| 2000 |         10 M （网页） | 100 MB |   1 GF (Intel Core) |
| 2010 |         10 G （广告） |   1 GB | 1 TF (Nvidia C2050) |
| 2020 |      1 T （社交网络） | 100 GB | 1 PF (Nvidia DGX-2) |

很明显，随机存取存储器没有跟上数据增长的步伐。 与此同时，算力的增长速度已经超过了现有数据的增长速度。 这意味着统计模型需要提高内存效率（这通常是通过添加非线性来实现的），同时由于计算预算的增加，能够花费更多时间来优化这些参数。 因此，机器学习和统计的关注点从（广义的）线性模型和核方法转移到了深度神经网络。 这也造就了许多深度学习的中流砥柱，如多层感知机 ([McCulloch and Pitts, 1943]
(https://zh-v2.d2l.ai/chapter\_references/zreferences.htmlid106)) 
、卷积神经网络 ([LeCun et al., 1998]
(https://zh-v2.d2l.ai/chapter\_references/zreferences.htmlid90)) 
、长短期记忆网络 ([Graves and Schmidhuber, 2005]
(https://zh-v2.d2l.ai/chapter\_references/zreferences.htmlid51)) 和Q学习 ([Watkins and Dayan, 1992](https://zh-v2.d2l.ai/chapter\_references/zreferences.htmlid179)) ，在相对休眠了相当长一段时间之后，在过去十年中被“重新发现”。

最近十年，在统计模型、应用和算法方面的进展就像寒武纪大爆发——历史上物种飞速进化的时期。 事实上，最先进的技术不仅仅是将可用资源应用于几十年前的算法的结果。 下面列举了帮助研究人员在过去十年中取得巨大进步的想法（虽然只触及了皮毛）。

- 新的容量控制方法，如dropout ([Srivastava et al., 2014], %(https://zh-v2.d2l.ai/chapter\_references/zreferences.htmlid155))，
有助于减轻过拟合的危险。这是通过在整个神经网络中应用噪声注入 ([Bishop, 1995](https://zh-v2.d2l.ai/chapter\_references/zreferences.htmlid9)) 来实现的，出于训练目的，用随机变量来代替权重。
- 注意力机制解决了困扰统计学一个多世纪的问题：如何在不增加可学习参数的情况下增加系统的记忆和复杂性。研究人员通过使用只能被视为可学习的指针结构 ([Bahdanau et al., 2014](https://zh-v2.d2l.ai/chapter\_references/zreferences.htmlid6)) 找到了一个优雅的解决方案。不需要记住整个文本序列（例如用于固定维度表示中的机器翻译），所有需要存储的都是指向翻译过程的中间状态的指针。这大大提高了长序列的准确性，因为模型在开始生成新序列之前不再需要记住整个序列。
- 多阶段设计。例如，存储器网络 ([Sukhbaatar et al., 2015]%(https://zh-v2.d2l.ai/chapter\_references/zreferences.htmlid158)) 
和神经编程器-解释器 ([Reed and De Freitas, 2015]。%(https://zh-v2.d2l.ai/chapter\_references/zreferences.htmlid136))。
它们允许统计建模者描述用于推理的迭代方法。这些工具允许重复修改深度神经网络的内部状态，从而执行推理链中的后续步骤，类似于处理器如何修改用于计算的存储器。
- 另一个关键的发展是生成对抗网络 ([Goodfellow et al., 2014]
%(https://zh-v2.d2l.ai/chapter\_references/zreferences.htmlid48)) 
的发明。传统模型中，密度估计和生成模型的统计方法侧重于找到合适的概率分布（通常是近似的）和抽样算法。
因此，这些算法在很大程度上受到统计模型固有灵活性的限制。生成式对抗性网络的关键创新是用具有可微参数的任意算法代替采样器。然后对这些数据进行调整，使得鉴别器（实际上是一个双样本测试）不能区分假数据和真实数据。通过使用任意算法生成数据的能力，它为各种技术打开了密度估计的大门。驰骋的斑马 ([Zhu et al., 2017]
%(https://zh-v2.d2l.ai/chapter\_references/zreferences.htmlid199)) 
和假名人脸 ([Karras et al., 2017]%(https://zh-v2.d2l.ai/chapter\_references/zreferences.htmlid81)) 
的例子都证明了这一进展。即使是业余的涂鸦者也可以根据描述场景布局的草图生成照片级真实图像（ ([Park et al., 2019]%(https://zh-v2.d2l.ai/chapter\_references/zreferences.htmlid120)) ）。
- 在许多情况下，单个GPU不足以处理可用于训练的大量数据。在过去的十年中，构建并行和分布式训练算法的能力有了显著提高。设计可伸缩算法的关键挑战之一是深度学习优化的主力——随机梯度下降，它依赖于相对较小的小批量数据来处理。同时，小批量限制了GPU的效率。因此，在1024个GPU上进行训练，例如每批32个图像的小批量大小相当于总计约32000个图像的小批量。最近的工作，首先是由 ([Li, 2017](https://zh-v2.d2l.ai/chapter\_references/zreferences.htmlid91)) 完成的，随后是 ([You et al., 2017](https://zh-v2.d2l.ai/chapter\_references/zreferences.htmlid192)) 和 ([Jia et al., 2018](https://zh-v2.d2l.ai/chapter\_references/zreferences.htmlid79)) ，将观察大小提高到64000个，将ResNet-50模型在ImageNet数据集上的训练时间减少到不到7分钟。作为比较——最初的训练时间是按天为单位的。
- 并行计算的能力也对强化学习的进步做出了相当关键的贡献。这导致了计算机在围棋、雅达里游戏、星际争霸和物理模拟（例如，使用MuJoCo）中实现超人性能的重大进步。有关如何在AlphaGo中实现这一点的说明，请参见如 ([Silver et al., 2016]
%(https://zh-v2.d2l.ai/chapter\_references/zreferences.htmlid152)) 。
简而言之，如果有大量的（状态、动作、奖励）三元组可用，即只要有可能尝试很多东西来了解它们之间的关系，强化学习就会发挥最好的作用。仿真提供了这样一条途径。
- 深度学习框架在传播思想方面发挥了至关重要的作用。允许轻松建模的第一代框架包括[Caffe]%(https://github.com/BVLC/caffe)
、[Torch]%(https://github.com/torch)
和[Theano]。%(https://github.com/Theano/Theano)。
许多开创性的论文都是用这些工具写的。到目前为止，它们已经被[TensorFlow]%(https://github.com/tensorflow/tensorflow)
（通常通过其高级API [Keras](https://github.com/keras-team/keras)使用）、[CNTK]%(https://github.com/Microsoft/CNTK)
、[Caffe 2]%(https://github.com/caffe2/caffe2)
和[Apache MXNet]%(https://github.com/apache/incubator-mxnet)
所取代。第三代工具，即用于深度学习的命令式工具，可以说是由[Chainer]
%(https://github.com/chainer/chainer)
率先推出的，它使用类似于Python NumPy的语法来描述模型。这个想法被[PyTorch]
%(https://github.com/pytorch/pytorch)
、MXNet的[Gluon API]%(https://github.com/apache/incubator-mxnet)
和[Jax]%(https://github.com/google/jax)
都采纳了。

“系统研究人员构建更好的工具”和“统计建模人员构建更好的神经网络”之间的分工大大简化了工作。 例如，在2014年，对卡内基梅隆大学机器学习博士生来说，训练线性回归模型曾经是一个不容易的作业问题。 而现在，这项任务只需不到10行代码就能完成，这让每个程序员轻易掌握了它。

 1.6. 深度学习的成功案例

人工智能在交付结果方面有着悠久的历史，它能带来用其他方法很难实现的结果。例如，使用光学字符识别的邮件分拣系统从20世纪90年代开始部署，毕竟，这是著名的手写数字MNIST数据集的来源。这同样适用于阅读银行存款支票和对申请者的信用进行评分。系统会自动检查金融交易是否存在欺诈。这成为许多电子商务支付系统的支柱，如PayPal、Stripe、支付宝、微信、苹果、Visa和万事达卡。国际象棋的计算机程序已经竞争了几十年。机器学习在互联网上提供搜索、推荐、个性化和排名。换句话说，机器学习是无处不在的，尽管它经常隐藏在视线之外。

直到最近，人工智能才成为人们关注的焦点，主要是因为解决了以前被认为难以解决的问题，这些问题与消费者直接相关。许多这样的进步都归功于深度学习。

- 智能助理，如苹果的Siri、亚马逊的Alexa和谷歌助手，都能够相当准确地回答口头问题。这包括一些琐碎的工作，比如打开电灯开关（对残疾人来说是个福音）甚至预约理发师和提供电话支持对话。这可能是人工智能正在影响我们生活的最明显的迹象。
- 数字助理的一个关键要素是准确识别语音的能力。逐渐地，在某些应用中，此类系统的准确性已经提高到与人类同等水平的程度 ([Xiong et al., 2018](https://zh-v2.d2l.ai/chapter\_references/zreferences.htmlid190))。
- 物体识别同样也取得了长足的进步。估计图片中的物体在2010年是一项相当具有挑战性的任务。在ImageNet基准上，来自NEC实验室和伊利诺伊大学香槟分校的研究人员获得了28%的Top-5错误率 ([Lin et al., 2010](https://zh-v2.d2l.ai/chapter\_references/zreferences.htmlid96)) 。到2017年，这一错误率降低到2.25% ([Hu et al., 2018](https://zh-v2.d2l.ai/chapter\_references/zreferences.htmlid72)) 。同样，在鉴别鸟类或诊断皮肤癌方面也取得了惊人的成果。
- 游戏曾经是人类智慧的堡垒。从TD-Gammon开始，一个使用时差强化学习的五子棋游戏程序，算法和计算的进步导致了算法被广泛应用。与五子棋不同的是，国际象棋有一个复杂得多的状态空间和一组动作。深蓝公司利用大规模并行性、专用硬件和高效搜索游戏树 ([Campbell et al., 2002](https://zh-v2.d2l.ai/chapter\_references/zreferences.htmlid19)) 击败了加里·卡斯帕罗夫(Garry Kasparov)。围棋由于其巨大的状态空间，难度更大。AlphaGo在2015年达到了相当于人类的棋力，使用和蒙特卡洛树抽样 ([Silver et al., 2016](https://zh-v2.d2l.ai/chapter\_references/zreferences.htmlid152)) 相结合的深度学习。扑克中的挑战是状态空间很大，而且没有完全观察到（我们不知道对手的牌）。在扑克游戏中，库图斯使用有效的结构化策略超过了人类的表现 ([Brown and Sandholm, 2017](https://zh-v2.d2l.ai/chapter\_references/zreferences.htmlid18)) 。这说明了游戏取得了令人瞩目的进步以及先进的算法在其中发挥了关键作用的事实。
- 人工智能进步的另一个迹象是自动驾驶汽车和卡车的出现。虽然完全自主还没有完全触手可及，但在这个方向上已经取得了很好的进展，特斯拉（Tesla）、英伟达（NVIDIA）和Waymo等公司的产品至少实现了部分自主。让完全自主如此具有挑战性的是，正确的驾驶需要感知、推理和将规则纳入系统的能力。目前，深度学习主要应用于这些问题的计算机视觉方面。其余部分则由工程师进行大量调整。

同样，上面的列表仅仅触及了机器学习对实际应用的影响之处的皮毛。 例如，机器人学、物流、计算生物学、粒子物理学和天文学最近取得的一些突破性进展至少部分归功于机器学习。 因此，机器学习正在成为工程师和科学家必备的工具。

关于人工智能的非技术性文章中，经常提到人工智能奇点的问题：机器学习系统会变得有知觉，并独立于主人来决定那些直接影响人类生计的事情。 在某种程度上，人工智能已经直接影响到人类的生计：信誉度的自动评估，车辆的自动驾驶，保释决定的自动准予等等。 甚至，我们可以让Alexa打开咖啡机。

幸运的是，我们离一个能够控制人类创造者的有知觉的人工智能系统还很远。 首先，人工智能系统是以一种特定的、面向目标的方式设计、训练和部署的。 虽然他们的行为可能会给人一种通用智能的错觉，但设计的基础是规则、启发式和统计模型的结合。 其次，目前还不存在能够自我改进、自我推理、能够在试图解决一般任务的同时，修改、扩展和改进自己的架构的“人工通用智能”工具。

一个更紧迫的问题是人工智能在日常生活中的应用。 卡车司机和店员完成的许多琐碎的工作很可能也将是自动化的。 农业机器人可能会降低有机农业的成本，它们也将使收割作业自动化。 工业革命的这一阶段可能对社会的大部分地区产生深远的影响，因为卡车司机和店员是许多国家最常见的工作之一。 此外，如果不加注意地应用统计模型，可能会导致种族、性别或年龄偏见，如果自动驱动相应的决策，则会引起对程序公平性的合理关注。 重要的是要确保小心使用这些算法。 就我们今天所知，这比恶意超级智能毁灭人类的风险更令人担忧。

 1.7. 特点

到目前为止，本节已经广泛地讨论了机器学习，它既是人工智能的一个分支，也是人工智能的一种方法。 虽然深度学习是机器学习的一个子集，但令人眼花缭乱的算法和应用程序集让人很难评估深度学习的具体成分是什么。 这就像试图确定披萨所需的配料一样困难，因为几乎每种成分都是可以替代的。

如前所述，机器学习可以使用数据来学习输入和输出之间的转换，例如在语音识别中将音频转换为文本。 在这样做时，通常需要以适合算法的方式表示数据，以便将这种表示转换为输出。 深度学习是“深度”的，模型学习了许多“层”的转换，每一层提供一个层次的表示。 例如，靠近输入的层可以表示数据的低级细节，而接近分类输出的层可以表示用于区分的更抽象的概念。 由于表示学习（representation learning）目的是寻找表示本身，因此深度学习可以称为“多级表示学习”。

本节到目前为止讨论的问题，例如从原始音频信号中学习，图像的原始像素值，或者任意长度的句子与外语中的对应句子之间的映射，都是深度学习优于传统机器学习方法的问题。 事实证明，这些多层模型能够以以前的工具所不能的方式处理低级的感知数据。 毋庸置疑，深度学习方法中最显著的共同点是使用端到端训练。 也就是说，与其基于单独调整的组件组装系统，不如构建系统，然后联合调整它们的性能。 例如，在计算机视觉中，科学家们习惯于将特征工程的过程与建立机器学习模型的过程分开。 Canny边缘检测器 ([Canny, 1987](https://zh-v2.d2l.ai/chapter\_references/zreferences.htmlid20)) 和SIFT特征提取器 ([Lowe, 2004](https://zh-v2.d2l.ai/chapter\_references/zreferences.htmlid102)) 作为将图像映射到特征向量的算法，在过去的十年里占据了至高无上的地位。 在过去的日子里，将机器学习应用于这些问题的关键部分是提出人工设计的特征工程方法，将数据转换为某种适合于浅层模型的形式。 然而，与一个算法自动执行的数百万个选择相比，人类通过特征工程所能完成的事情很少。 当深度学习开始时，这些特征抽取器被自动调整的滤波器所取代，产生了更高的精确度。

因此，深度学习的一个关键优势是它不仅取代了传统学习管道末端的浅层模型，而且还取代了劳动密集型的特征工程过程。 此外，通过取代大部分特定领域的预处理，深度学习消除了以前分隔计算机视觉、语音识别、自然语言处理、医学信息学和其他应用领域的许多界限，为解决各种问题提供了一套统一的工具。

除了端到端的训练，人们正在经历从参数统计描述到完全非参数模型的转变。 当数据稀缺时，人们需要依靠简化对现实的假设来获得有用的模型。 当数据丰富时，可以用更准确地拟合实际情况的非参数模型来代替。 在某种程度上，这反映了物理学在上个世纪中叶随着计算机的出现所经历的进步。 现在人们可以借助于相关偏微分方程的数值模拟，而不是用手来求解电子行为的参数近似。这导致了更精确的模型，尽管常常以牺牲可解释性为代价。

与以前工作的另一个不同之处是接受次优解，处理非凸非线性优化问题，并且愿意在证明之前尝试。 这种在处理统计问题上新发现的经验主义，加上人才的迅速涌入，导致了实用算法的快速进步。 尽管在许多情况下，这是以修改和重新发明存在了数十年的工具为代价的。

最后，深度学习社区引以为豪的是，他们跨越学术界和企业界共享工具，发布了许多优秀的算法库、统计模型和经过训练的开源神经网络。 正是本着这种精神，本书免费分发和使用。我们努力降低每个人了解深度学习的门槛，希望读者能从中受益。

 1.8. 小结

- 机器学习研究计算机系统如何利用经验（通常是数据）来提高特定任务的性能。它结合了统计学、数据挖掘和优化的思想。通常，它是被用作实现人工智能解决方案的一种手段。
- 表示学习作为机器学习的一类，其研究的重点是如何自动找到合适的数据表示方式。深度学习是通过学习多层次的转换来进行的多层次的表示学习。
- 深度学习不仅取代了传统机器学习的浅层模型，而且取代了劳动密集型的特征工程。
- 最近在深度学习方面取得的许多进展，大都是由廉价传感器和互联网规模应用所产生的大量数据，以及（通过GPU）算力的突破来触发的。
- 整个系统优化是获得高性能的关键环节。有效的深度学习框架的开源使得这一点的设计和实现变得非常容易。

 1.9. 练习

1. 你当前正在编写的代码的哪些部分可以“学习”，即通过学习和自动确定代码中所做的设计选择来改进？你的代码是否包含启发式设计选择？
2. 你遇到的哪些问题有许多解决它们的样本，但没有具体的自动化方法？这些可能是使用深度学习的主要候选者。
3. 如果把人工智能的发展看作一场新的工业革命，那么算法和数据之间的关系是什么？它类似于蒸汽机和煤吗？根本区别是什么？
4. 你还可以在哪里应用端到端的训练方法，比如 [图1.1.2](https://zh-v2.d2l.ai/chapter\_introduction/index.htmlfig-ml-loop) 、物理、工程和计量经济学？



<iframe src="https://discuss.d2l.ai/embed/comments?topic\_id=1744" id="discourse-embed-frame" width="100%" frameborder="0" scrolling="no" referrerpolicy="no-referrer-when-downgrade" height="8170px" style="vertical-align: middle; display: block; width: 706px; border: none;"></iframe>



这本书写于一个翻天覆地的时代，技术的演变让我们重新审视贝叶斯公式以及它在知识大厦中的位置。 

这本书也写在了一个传播方式改变了我们谈论科学方式的时代。



贝叶斯主义就是假设“现实”的所有模型、理 论或概念都只不过是某种信念、虚构或诗歌，尤其要指出的是，“所有 模型都是错的”;然后，实际数据应该迫使我们调整赋予不同模型的重 要性，即置信度;关键在于，调整这些置信度的方式应该尽可能严谨 地遵循贝叶斯公式。

```
我之前一直觉得这种知识哲学没有意义。它似乎否定了所有关于现实
和真理的概念，即使这些概念对研究人员来说如此重要;但它又似乎
```

完全符合物理学家、诺贝尔物理学奖获得者理查德·费曼的说法 [4] : “我能带着疑问、不确定和无知活着。我觉得，比起知道一些可能错误 的答案，还是不知道答案的生活更有趣。我有些近似的答案，对于各 种问题也有些确定程度或高或低的合理信念，但我不会绝对确信任何 事情。也有很多我一点都不明白的事情，但我不一定需要一个答案。 我不害怕‘我不知道’这个事实。”

你可能会觉得这种观点激动人心，或者希望将这种看待知识的方法拒
之门外。然而在选择接受还是拒绝贝叶斯主义之前，我只能鼓励你先
花点时间，长考一下贝叶斯公式与它的推论。

遗憾的是，这本书里的主要向导，也就是我，对这个公式的理解还非 常浅薄。



工程思维/应用思维：



正是这种限制迫使我们只能考虑一种必定仅作为近似的贝叶斯主义， 我称之为实用贝叶斯主义。它跟纯粹贝叶斯主义的不同之处，就是它 要求(迅速的)可计算性。

可能你也猜到了，理解纯粹贝叶斯主义者与实用贝叶斯主义者并不容 易。为此，我们必须考虑大量的基础概念，它们来自数学、逻辑学、 统计学、计算机科学、人工智能，甚至还有物理学、生物学、神经科 学、心理学和经济学。我们需要谈到对数、逆否命题、 值、所罗门 诺夫复杂度和神经网络，还有熵、达尔文式演化、虚假回忆、认知偏 差和金融泡沫。另外，我们也会用到科学史中的大量例子来考验我们 这两位虚构人物。

对，我知道为了理解贝叶斯公式，要明白的东西可不少......

但这岂不是正好，因为我喜欢在闲暇时解释现代科学，我甚至为此开 通了名为 Science4All 的 YouTube 频道!所以，与其把这本书当成哲 学书来读，我请大家不如把它当成科学和数学的科普书。另外，在解 释贝叶斯主义的途中，我也不惜绕几个弯，拐到科学的漫谈上，暗地 里就是为了鼓励你在了解科学理论的道路上走得更远!



人类实在太有智慧了，希望能越来越向她靠近.........即使在开始教授这门课以及撰写这本书之后，我仍然一遍又一遍地从这个主题中发现了不可胜数的惊 人奇景，它已经变成我最喜欢的数学等式了。





要以理性的方式谈论理性，就 要先用理性的方式定义理性......这就像一条咬着自己尾巴的蛇。

这个难点当然并非贝叶斯公式所独有，所有知识哲学似乎都必然受制
于这种自我指涉。数学家也曾花上数个世纪的努力来发展没有自我指
涉的理论，然而并不成功。(哥德尔，谢谢你!)



严谨、简洁与清晰。它定义了如此清晰 的推理规则 3，应用这些规则似乎足以相对精确地(即使只是近似 地)理解这个世界。这正是计算机科学家的理想，只要按下启动按 钮，机器就能执行一系列指令来自动达到目标。这说的当然就是人工 智能!30 年以来，贝叶斯公式一直处于这个领域中众多研究的核心， 这大概并非偶然。



法国的打分制以 20 分为满分。



即使指明用到了什么方法也不够。利用机器学习在大数据中推 断出有用信息的数据科学家很早就发现了，没有人工干扰不一定能保 证客观性。无论是人还是机器，我们似乎都必定要在某个模型内部进 行推理。所以说，我们的结论似乎必然依赖于模型。纯粹贝叶斯主义 者断定，这就说明了所有知识都必然是主观的。



可惜的是，这本书跟一切篇幅有限的书一样，难免挂一漏万。我在这
里先为本书不可胜数的不足之处说声抱歉，尤其是因为我没有花时间
将贝叶斯主义与其他与之竞争的哲学流派进行深入比较。我的目标没
那么远大，只是希望能帮助你理解贝叶斯主义的某些重要方面，至少
是我理解的那些方面。的确，跟本书一样，我的大脑也是有限的，所
以请你原谅我的无知之处。我会尝试提及和揭示尽可能多的对我来说
重要的思想，但也必定会遗漏那些我不知道或误判其重要性的内容。

另外，本书描述的是截至出版时我的认知状态，但在掌握贝叶斯主义 的旅程中，我希望之后能取得更多的进展。如果你愿意伴随我踏上这 段旅程的话，请在 Twitter 上关注我(@science4all)，并查看我的视 频频道 Science4All，我从 2018 年末开始在频道中放了一系列关于贝 叶斯公式的视频。我同样邀请你前往我和哲学家蒂博·吉罗(网名是 Monsieur Phi)共同主持的播客“公理”(Axiome)，倾听我们的思 考。尽管我们的目标是探讨所有与数学、哲学和各种自然科学有关的 东西，但鉴于我们都对概率和逻辑无比着迷，我们也反复谈到了贝叶 斯主义的方方面面，其中一些内容在本书中没有提及，比如贝特朗悖 论的贝叶斯解释 [12]。



读这本书的时候请不要急躁，要多花点时间思考，但也不要轻易放 弃。

读这本书的时候请不要急躁，要多花点时间思考，但也不要轻易放 弃。这本书并不是越往后面越难，你可以在没有阅读之前章节的情况 下享受每一章——虽然按顺序阅读每一章可能更好。这不是一本教材，也没有考试。我不会要求你理解所有内容，甚至强烈建议你跳过
那些太复杂的段落继续阅读(如果你之后肯回到这些难度较大的段落
的话)。我的目标不是让你成为贝叶斯主义的专家。

我最希望的是，你能享受贝叶斯的推理以及在理解贝叶斯主义的基础
和推论时用到的那些科学内容，并从中找到美感。我希望你能把自己
当作一位探险家，出发去探索未知的土地，发现各种各样的动物、植
物、文化与引人入胜的风景，而不一定要花时间学会本地语言中的所
有细微之处。我希望你能从这段旅程中获益。
如果你跟随我的脚步的话，那么我希望你也会沉浸在热忱、魅力和疑
问之中，这就是本书的首要目的。



学习是一支舞蹈。

正如一般的数学内容，这个公式的抽象性和复杂程度足以吓退我们之中不够勇敢的那些人。



想要做到这一点，，毅力与精进是不可或缺的。

要付出的代价就是大量脑力劳动，但其回报绝对丰厚。对于纯粹贝叶 斯主义者来说，最终能够(足够)正确思考的能力就在航程的终点。



贝叶斯公式的优雅，还有它的推论将我引向了不可胜数的快乐思考，
从中得出的知识哲学让我一直感受着快乐和幸福。正因如此，我最终
得出，贝叶斯公式是数学中最优美的等式。



对于纯粹贝叶斯主义者来说，用理性的方式研究历
史是完全可行的。对于某些人来说，历史、物种演化与宇宙学这些学
科并非科学，因为它们不能通过可重复的实验来研究。有趣的是，这
种思考不过是来自波普尔哲学以及频率主义统计学的单纯假象。



所有人类知识从直觉开始，然后转变为观念，最后化为思想。 伊曼努尔·康德(1724—1804)



根据贝叶斯定理，任何理论都不完美。取而代之的是一项未竟的
工作，它永远处于推敲与测试之中。





====
前言 数学踏上新的征程

人们常说，刚刚过去的一个世纪是数学真正的黄金时代。数学在 20 世 纪的进步比先前所有的世纪加起来还要大。然而，刚刚开始的这个世纪 也可能同样是数学发展的好时候。或许，数学在这个世纪的变迁会和 20 世纪一样巨大，甚至更为惊人。引发这种想法的信号之一是一场渐 变:自 20 世纪 70 年代开始，数学方法的基石——证明的概念逐渐发生 演变，让一个古老却有些被人忽视的数学概念重新回到了舞台中央，这 就是“计算”。

计算能成为引发革命的导火索，这看起来有点不合常理。算法，比如做 加法和做乘法的算法，常常被视为数学知识中最基础的一部分，做计算 也经常被看成是缺乏创造性的枯燥工作。数学家们自己对计算也颇有成 见，勒内·托姆就曾说过:“我的论述中很大一部分属于纯粹的猜想，大 家基本上可以把它们看成是梦话。我接受这种定性......如今，世界上到 处有这么多学者在做计算，难道有人做梦不是件好事吗?”用计算来做 梦，大概还真有点难度啊......

不幸的是，对计算的偏见恰恰根植于“数学证明”这一概念的定义里。确 实，欧几里得以降，“证明”的定义就是利用公理和演绎规则构建的一套 推理 。然而，要解决一个数学问题，仅仅需要构建一套推理吗?数学 的实践难道没有告诉我们，解决问题需要把推理的步骤和计算的步骤巧 妙地融合起来吗?公理化方法若局限在推理中，它所展现的数学视野恐 怕也会十分狭隘。正是因为人们对约束过多的公理化方法多有批评，才 让计算有机会重新出现在数学的舞台上。现在，已有一些研究工作(它 们之间未必有关联)渐渐开始质疑推理高于计算的优势地位，并倡导一 种更为平衡的观点，让两者互为补充。

这场革命让我们重新考量推理和计算之间的关系，同时也促使我们重新 审视数学与物理学、生物学等自然科学之间的对话，特别是数学为何能 在这些学科中发挥难以理解的强大作用这一古老问题，以及自然理论的 逻辑形式这一全新问题。此外，这场革命给“分析判断”和“综合判断”等 哲学概念带来了新的火花。它还让我们反思数学与计算机科学之间的关 系，而且数学似乎是唯一一门不需要借助机器的科学，它为什么如此独 特?

最后，最振奋人心的是，这场革命让我们隐约看到了一些解决数学问题 的新方式，它摆脱了过去的技术强加给证明长度的枷锁——数学也许正 踏上新的征程，去探索从未涉足的全新领域。

诚然，公理化方法的危机并不是凭空出现的。从 20 世纪上半叶起就有 许多先兆，特别是两种理论——可计算性理论和构造性理论的出现。这 两种理论本身虽然没有质疑公理化方法，却重新确立了计算在数学大厦 中的地位。在讨论公理化危机之前，我们会简要回顾这两个概念的历 史。不过，还是让我们先上溯远古，探寻计算这一概念的起源，看看古 希腊人对数学的“发明”过程吧。


目标读者
本书面向学生（本科生或研究生）、工程师和研究人员，他们希望扎实掌握深度学习的实用技术。因为我们从头开始解释每个概念，所以不需要过往的深度学习或机器学习背景。全面解释深度学习的方法需要一些数学和编程，但我们只假设读者了解一些基础知识，包括线性代数、微积分、概率和非常基础的Python编程。此外，在附录中，我们提供了本书所涵盖的大多数数学知识的复习。大多数时候，我们会优先考虑直觉和想法，而不是数学的严谨性。有许多很棒的书可以引导感兴趣的读者走得更远。BelaBollobas的《线性分析》([Bollobás,1999](https：//zh-v2.d2l.ai/chapter\_references/zreferences.htmlid13))对线性代数和函数分析进行了深入的研究。([Wasserman,2013](https：//zh-v2.d2l.ai/chapter\_references/zreferences.htmlid178))是一本很好的统计学指南。


-数学思维与人工智能的关系
-从数学到智能的发展历程


\newpage